% =====================================================================
%  Marés Oceânicas: da dinâmica newtoniana às Equações de Maré de Laplace
%  Documento de estudo
%  Compilar com:  latexmk -pdf mares-oceanicas.tex
% =====================================================================
\documentclass[11pt,a4paper,oneside]{report}

\usepackage[utf8]{inputenc}
\usepackage[T1]{fontenc}
\usepackage{newpxtext}
\usepackage{newpxmath}
\usepackage{inconsolata}
\usepackage{microtype}

\usepackage[a4paper,margin=2.6cm,bottom=3.0cm,headheight=14pt]{geometry}
\let\Bbbk\relax
\let\openbox\relax
\usepackage{amsmath,amssymb,amsthm,mathtools}
\usepackage{empheq}
\usepackage{bm}
\usepackage{graphicx}
\usepackage{booktabs,array,multirow}
\usepackage{siunitx}
\usepackage{xcolor}
\usepackage{tikz}
\usepackage{pgfplots}
\usepackage{enumitem}
\usepackage{caption}
\usepackage{subcaption}
\usepackage{float}
\usepackage{listings}
\usepackage{titlesec}
\usepackage{fancyhdr}
\usepackage[breakable,skins]{tcolorbox}
\usepackage[hidelinks,colorlinks=true,linkcolor=azulforte,citecolor=azulforte,
            urlcolor=azulforte,
            pdftitle={Mares Oceanicas: da dinamica newtoniana as Equacoes de Mare de Laplace},
            pdfsubject={Marés oceânicas: formulação e modelagem numérica},
            pdfkeywords={mares, Laplace, dinamica classica, analise harmonica,
                         pontos anfidromicos, integradores simpleticos}]{hyperref}

\usetikzlibrary{arrows.meta,decorations.pathmorphing,patterns,calc,positioning}
\pgfplotsset{compat=1.18}
\graphicspath{{../../media/}{media/}}

% --- português sem babel (o texlive local não traz babel-portuguese) ---
\renewcommand{\chaptername}{Capítulo}
\renewcommand{\partname}{Parte}
\renewcommand{\contentsname}{Sumário}
\renewcommand{\appendixname}{Apêndice}
\renewcommand{\bibname}{Referências}
\renewcommand{\figurename}{Figura}
\renewcommand{\tablename}{Tabela}
\renewcommand{\abstractname}{Resumo}
\renewcommand{\listfigurename}{Lista de figuras}
\renewcommand{\listtablename}{Lista de tabelas}

% padrões de hifenização são os do inglês; afrouxa o ajuste e corrige à mão
% as palavras técnicas que mais se repetem
\tolerance=1400
\emergencystretch=2.4em
\hyphenation{an-fi-dro-mo an-fi-dro-mos an-fi-dro-mi-co an-fi-dro-mi-cos
  equi-li-brio hi-dro-di-na-mi-ca hi-dro-di-na-mi-co cons-ti-tuin-te
  cons-ti-tuin-tes ba-ti-me-tria dis-si-pa-cao ba-ro-tro-pi-co
  ba-ro-cli-ni-co sim-ple-ti-co sim-ple-ti-cos har-mo-ni-ca har-mo-ni-cas
  ma-re-gra-fo ma-re-gra-fos os-ci-la-dor res-so-nan-cia gra-vi-ta-cio-nal
  per-tur-ba-dor La-pla-ce New-ton amor-te-ci-men-to pro-fun-di-da-de
  se-mi-diur-na se-mi-diur-nas nu-me-ri-ca nu-me-ri-co geo-po-ten-cial}

% --- cores ---
\definecolor{azulforte}{HTML}{184F95}
\definecolor{azulclaro}{HTML}{EAF3FD}
\definecolor{tinta}{HTML}{1A1A1A}
\definecolor{cinza}{HTML}{5A5A5A}
\definecolor{laranja}{HTML}{EB6834}
\definecolor{verde}{HTML}{1BAF7A}
\definecolor{papel}{HTML}{FBFAF7}

% --- cabeçalhos ---
\pagestyle{fancy}
\fancyhf{}
\renewcommand{\headrulewidth}{0.4pt}
\fancyhead[L]{\small\scshape\nouppercase{\leftmark}}
\fancyhead[R]{\small\thepage}
\fancypagestyle{plain}{\fancyhf{}\renewcommand{\headrulewidth}{0pt}%
  \fancyfoot[C]{\small\thepage}}

\titleformat{\chapter}[display]
  {\normalfont\huge\bfseries\color{azulforte}}
  {\normalsize\scshape\color{cinza}\chaptertitlename\ \thechapter}{8pt}{\Huge}
\titlespacing*{\chapter}{0pt}{-10pt}{28pt}
\titleformat{\section}{\normalfont\Large\bfseries\color{azulforte}}{\thesection}{0.7em}{}
\titleformat{\subsection}{\normalfont\large\bfseries}{\thesubsection}{0.7em}{}

% --- ambientes ---
\theoremstyle{definition}
\newtheorem{definicao}{Definição}[chapter]
\newtheorem{exemplo}[definicao]{Exemplo}
\theoremstyle{plain}
\newtheorem{proposicao}[definicao]{Proposição}
\theoremstyle{remark}
\newtheorem{observacao}[definicao]{Observação}

\newtcolorbox{caixa}[1][]{
  breakable, enhanced, colback=azulclaro, colframe=azulforte,
  boxrule=0.4pt, left=8pt, right=8pt, top=6pt, bottom=6pt,
  fonttitle=\bfseries\small, title={#1}, sharp corners=downhill,
  attach boxed title to top left={xshift=8pt,yshift=-8pt},
  boxed title style={colback=azulforte,sharp corners}}

\newtcolorbox{alerta}[1][]{
  breakable, enhanced, colback=papel, colframe=laranja,
  boxrule=0.9pt, left=8pt, right=8pt, top=6pt, bottom=6pt,
  fonttitle=\bfseries\small, title={#1}, sharp corners,
  attach boxed title to top left={xshift=8pt,yshift=-8pt},
  boxed title style={colback=laranja,sharp corners}}

% --- código ---
\lstdefinestyle{py}{
  language=Python, basicstyle=\ttfamily\footnotesize,
  keywordstyle=\color{azulforte}\bfseries, commentstyle=\color{cinza},
  stringstyle=\color{verde}, numbers=left, numberstyle=\tiny\color{cinza},
  numbersep=8pt, frame=leftline, framerule=1.2pt, rulecolor=\color{azulclaro},
  backgroundcolor=\color{papel}, showstringspaces=false, breaklines=true,
  columns=fullflexible, xleftmargin=14pt, extendedchars=true, upquote=true,
  literate=%
    {á}{{\'a}}1 {à}{{\`a}}1 {ã}{{\~a}}1 {â}{{\^a}}1
    {é}{{\'e}}1 {ê}{{\^e}}1 {í}{{\'i}}1 {ó}{{\'o}}1
    {ô}{{\^o}}1 {õ}{{\~o}}1 {ú}{{\'u}}1 {ç}{{\c c}}1
    {Á}{{\'A}}1 {Ã}{{\~A}}1 {É}{{\'E}}1 {Ê}{{\^E}}1
    {Í}{{\'I}}1 {Ó}{{\'O}}1 {Ô}{{\^O}}1 {Õ}{{\~O}}1 {Ú}{{\'U}}1 {Ç}{{\c C}}1
}
\lstset{style=py}

% --- macros ---
\newcommand{\vect}[1]{\bm{#1}}
\newcommand{\uvec}[1]{\hat{\bm{#1}}}
\newcommand{\pd}[2]{\frac{\partial #1}{\partial #2}}
\newcommand{\ddt}[1]{\frac{\mathrm{d} #1}{\mathrm{d} t}}
\newcommand{\dd}{\mathrm{d}}
\newcommand{\ii}{\mathrm{i}}
\newcommand{\ordem}[1]{\mathcal{O}\!\left(#1\right)}
\DeclareMathOperator{\atantwo}{atan2}
\DeclareMathOperator{\sgn}{sgn}
\DeclareSIUnit{\ano}{ano}
\DeclareSIUnit{\anos}{anos}
\sisetup{output-decimal-marker={,},per-mode=symbol,group-separator={\,}}

\setlist[itemize]{itemsep=2pt,topsep=4pt}
\setlist[enumerate]{itemsep=2pt,topsep=4pt}

% =====================================================================
\begin{document}

\begin{titlepage}
\centering
\vspace*{2.4cm}
{\scshape\large Documento de estudo\par}
\vspace{0.4cm}
{\color{azulforte}\rule{\textwidth}{1.2pt}}
\vspace{0.9cm}

{\Huge\bfseries Marés Oceânicas\par}
\vspace{0.6cm}
{\LARGE Da dinâmica newtoniana às\\[2pt] Equações de Maré de Laplace\par}
\vspace{0.9cm}
{\color{azulforte}\rule{\textwidth}{1.2pt}}
\vspace{1.4cm}

{\large Fundamentos de mecânica clássica, formulação do problema de maré\\
e a matemática das melhorias implementadas na simulação\par}

\vfill
{\normalsize Texto de referência para a formulação newtoniana:\\
J.~B. Marion \& S.~T. Thornton, \emph{Classical Dynamics of Particles and\\
Systems}, \S5.5 ``Ocean Tides''\par}
\vspace{1.4cm}
{\large Agosto de 2026\par}
\end{titlepage}

\tableofcontents
\cleardoublepage

\chapter*{Apresentação}
\addcontentsline{toc}{chapter}{Apresentação}
\markboth{APRESENTAÇÃO}{}

O objeto deste texto é a maré oceânica, tratada do princípio ao fim: da
formulação newtoniana da força de maré num referencial acelerado até um modelo
barotrópico global das Equações de Maré de Laplace, confrontado com marégrafos.
O leitor pressuposto é o de pós-graduação em física, com mecânica clássica
madura e sem familiaridade prévia com oceanografia dinâmica.

A escolha de escopo é deliberada. A teoria de equilíbrio de Newton, que é o que
os livros de mecânica clássica apresentam, está \emph{qualitativamente} errada —
não apenas imprecisa —, e mostrar por quê exige atravessar hidrodinâmica de
fluido girante, dissipação turbulenta, gravitação autoconsistente e análise
harmônica. Esse percurso é longo o bastante para que o encadeamento se perca se
não for explicitado, e por isso o texto o expõe em quatro partes com uma
progressão única.

\paragraph{Método.} Três compromissos organizam a exposição.

\begin{enumerate}
\item \textbf{Toda equação é derivada, e toda derivação parte da anterior.}
  Resultados importados da literatura aparecem com a referência ao lado e são,
  sempre que possível, verificados por consistência interna — a força expandida
  contra a exata, o potencial contra o campo de forças que dele deriva, a
  resposta integrada contra a admitância analítica.
\item \textbf{Todo número citado é reproduzível.} As tabelas de resultados são
  saídas de execução do modelo descrito, não estimativas de ordem de grandeza,
  e as constantes usadas estão declaradas.
\item \textbf{Os limites de cada modelo são parte do conteúdo.} Há um capítulo
  inteiro sobre onde a teoria de equilíbrio falha, e seções sobre onde o modelo
  dinâmico erra e por quê. Um modelo cujo domínio de validade não foi
  delimitado não foi compreendido.
\end{enumerate}

\paragraph{Organização.} A Parte~\ref{parte:fundamentos} recolhe os fundamentos
de dinâmica clássica de que tudo o mais depende: referenciais não inerciais,
expansão multipolar do potencial gravitacional, oscilador forçado amortecido e o
mínimo necessário de dinâmica de fluidos geofísicos. Quem domina esse material
pode lê-la em diagonal, mas as convenções de sinal ali fixadas são usadas até o
fim. A Parte~\ref{parte:mare} aplica essa dinâmica ao problema de maré e chega à
teoria de equilíbrio, com os resultados clássicos reproduzidos numericamente. A
Parte~\ref{parte:melhorias} é o núcleo do texto: o fundamento teórico e a
formulação matemática do modelo dinâmico — Equações de Maré de Laplace,
auto-atração e carga, dissipação, pontos anfidrômicos, análise harmônica e
esquemas numéricos. A Parte~\ref{parte:sintese} fecha com validação contra
observação, orçamento energético e a distância que separa este modelo dos
operacionais.

\paragraph{Percursos de leitura.} Para o argumento essencial em cinco leituras:
\S\ref{sec:multipolo}, Capítulo~\ref{cap:forcamare}, \S\ref{sec:numeros},
Capítulo~\ref{cap:falha} e \S\ref{sec:topologia}. Quem se interessa
especificamente por análise de séries de marégrafo encontra no
Capítulo~\ref{cap:harmonica} um tratamento autocontido, que não depende da
hidrodinâmica. O Apêndice~\ref{ap:python} descreve, de forma sucinta, a
implementação numérica que produziu os resultados apresentados.




% =====================================================================
\part{Fundamentos de dinâmica clássica}
\label{parte:fundamentos}
% =====================================================================

\chapter{Referenciais não inerciais}
\label{cap:referenciais}

A maré é, antes de mais nada, um efeito de referencial. Um observador em queda
livre no campo de outro corpo não sente a gravidade desse corpo — sente apenas
a \emph{variação} dela através da sua própria extensão. É por isso que a maré
existe, e é por isso que a primeira coisa a fixar é a mecânica do referencial
não inercial.

\section{O operador de derivada temporal em referencial girante}

Sejam $S$ um referencial inercial e $S'$ um referencial cuja origem se move e
cujos eixos giram com velocidade angular $\vect{\omega}(t)$ em relação a $S$.
Para qualquer vetor $\vect{A}$,
\begin{equation}
\left(\ddt{\vect{A}}\right)_{S} \;=\; \left(\ddt{\vect{A}}\right)_{S'} \;+\;
\vect{\omega}\times\vect{A}.
\label{eq:operador}
\end{equation}
A demonstração é direta: escreva $\vect{A}=\sum_i A_i\,\uvec{e}_i'$ na base de
$S'$; a derivada em $S$ atinge tanto as componentes $A_i$ (o que dá o primeiro
termo) quanto os versores $\uvec{e}_i'$, cuja taxa de variação, para uma base
rígida em rotação, é $\dot{\uvec{e}}_i'=\vect{\omega}\times\uvec{e}_i'$.

A Eq.~\eqref{eq:operador} é um \emph{operador}: vale para qualquer vetor, e em
particular pode ser aplicada duas vezes.

\section{A equação de movimento no referencial acelerado e girante}

Seja $\vect{R}_0(t)$ a posição da origem de $S'$ vista de $S$, e $\vect{r}$ a
posição de uma partícula medida em $S'$. Então
$\vect{r}_{S}=\vect{R}_0+\vect{r}$, e aplicando \eqref{eq:operador} duas vezes
obtém-se a forma canônica
\begin{equation}
m\,\ddot{\vect{r}}\Big|_{S'} \;=\;
\underbrace{\vect{F}}_{\text{forças reais}}
\;-\;\underbrace{m\ddot{\vect{R}}_0}_{\text{arrasto de translação}}
\;-\;\underbrace{m\,\dot{\vect{\omega}}\times\vect{r}}_{\text{Euler}}
\;-\;\underbrace{2m\,\vect{\omega}\times\dot{\vect{r}}}_{\text{Coriolis}}
\;-\;\underbrace{m\,\vect{\omega}\times(\vect{\omega}\times\vect{r})}_{\text{centrífuga}}.
\label{eq:naoinercial}
\end{equation}

Os cinco termos aparecem todos neste documento, e é útil já dizer onde:

\begin{itemize}
\item $\vect{F}$ contém a gravidade da própria Terra e a atração do corpo
  perturbador (Lua, Sol);
\item $-m\ddot{\vect{R}}_0$ é o termo que \textbf{gera a maré}
  (Capítulo~\ref{cap:forcamare}): a origem escolhida é o centro da Terra, que
  está em queda livre no campo da Lua;
\item o termo de Euler é desprezível aqui ($\dot{\vect\Omega}_{\text{Terra}}\sim
  6\times10^{-22}\,\si{\radian\per\second\squared}$, associado ao próprio atrito de
  maré — ver \S\ref{sec:torque});
\item Coriolis é o termo que produz os \textbf{pontos anfidrômicos}
  (Capítulo~\ref{cap:anfidromos}); no oceano ele \emph{não} é uma correção
  pequena, como se mostrará em \S\ref{sec:escala};
\item a centrífuga é absorvida na definição de \emph{geoide}: o que se chama
  ``vertical'' e o que se mede como $g=\SI{9.80665}{\meter\per\second\squared}$
  já incluem $-\vect{\Omega}\times(\vect{\Omega}\times\vect{r})$. Ela achata a
  Terra em $1/298$ e é estática — não gera maré.
\end{itemize}

\begin{caixa}[Por que a centrífuga não é maré e o arrasto de translação é]
Ambas são forças inerciais, mas com estruturas espaciais distintas. A
centrífuga associada à \emph{rotação da Terra} é estacionária no referencial
que gira com o planeta: deforma o geoide uma vez, permanentemente. Já
$-m\ddot{\vect{R}}_0$ é uniforme no espaço mas o campo real do perturbador não
é, e a maré é exatamente a diferença entre os dois. A força de maré é, em uma
frase, \emph{o resíduo de um campo não uniforme depois de subtraída a queda
livre do centro}.
\end{caixa}

\section{O referencial geocêntrico}
\label{sec:geocentrico}

Adota-se, do Capítulo~\ref{cap:forcamare} em diante, o referencial com origem
no centro de massa da Terra. Ele é não inercial por duas razões independentes:

\begin{enumerate}
\item \textbf{translada aceleradamente}, porque a Terra orbita o baricentro
  Terra--Lua (a \SI{4671}{\kilo\meter} do centro da Terra, ou seja, dentro do
  planeta) e o baricentro Terra--Sol;
\item \textbf{gira}, com
  $\Omega_\oplus=2\pi/T_{\text{sid}}=\SI{7.292115e-5}{\radian\per\second}$,
  $T_{\text{sid}}=\SI{86164.0905}{\second}$.
\end{enumerate}

O ponto conceitual mais importante do documento talvez seja este: os dois
efeitos entram em \emph{lugares diferentes} da física da maré. A translação
acelerada gera a \textbf{força}; a rotação governa a \textbf{resposta}. A
teoria de equilíbrio (Parte~\ref{parte:mare}) usa apenas a primeira; as
Equações de Maré de Laplace (Parte~\ref{parte:melhorias}) usam as duas, e é daí
que vem a diferença qualitativa entre os dois modelos.

\begin{observacao}
A rotação da Terra é praticamente uniforme na escala de tempo de interesse, de
modo que $\dot{\vect\Omega}\approx 0$ e o termo de Euler some da
Eq.~\eqref{eq:naoinercial}. Mas ele não é rigorosamente nulo, e a sua
não-nulidade \emph{é} um efeito de maré: o atrito das marés frena a rotação
terrestre em $\sim\SI{2.3}{\milli\second} $ por século
(\S\ref{sec:torque}).
\end{observacao}

% =====================================================================
\chapter{Gravitação, dois corpos e a expansão multipolar}
\label{cap:gravitacao}

\section{O problema de dois corpos}

Para duas massas $m_1,m_2$ interagindo apenas gravitacionalmente, a separação
$\vect{d}=\vect{r}_2-\vect{r}_1$ obedece
\begin{equation}
\mu\,\ddot{\vect{d}} = -\frac{G m_1 m_2}{d^{3}}\vect{d},
\qquad \mu=\frac{m_1m_2}{m_1+m_2},
\end{equation}
isto é, o problema de dois corpos reduz-se a \emph{um} corpo de massa reduzida
$\mu$ num campo central $-GM_{\text{tot}}/d$. A conservação do momento angular
$\vect{L}=\mu\,\vect{d}\times\dot{\vect{d}}$ confina o movimento a um plano, e
a conservação da energia leva à órbita cônica
\begin{equation}
d(\theta)=\frac{a(1-e^{2})}{1+e\cos\theta}.
\end{equation}

Para a Lua, $a=\SI{3.844e8}{\meter}$, $e=0.0549$, inclinação
$i=\ang{5.145}$ em relação à eclíptica; para a órbita da Terra em torno do Sol,
$a=\SI{1.4959787e11}{\meter}$, $e=0.016708$, e a obliquidade da eclíptica é
$\varepsilon=\ang{23.4393}$.

\begin{caixa}[Por que a excentricidade importa muito mais do que parece]
A maré escala com $D^{-3}$ (Capítulo~\ref{cap:potencial}). Logo
\begin{equation}
\frac{\delta h}{h} = -3\,\frac{\delta D}{D},
\label{eq:amplificacao3x}
\end{equation}
e qualquer erro de distância é \textbf{amplificado três vezes} na amplitude da
maré. Com $D$ variando de \SI{363625}{\kilo\meter} (perigeu) a
\SI{405870}{\kilo\meter} (apogeu), a modulação da amplitude por esse efeito
sozinho é $(405870/363625)^3=\num{1.39}$: as \emph{marés perigeanas}. A
Eq.~\eqref{eq:amplificacao3x} volta a aparecer no Capítulo~\ref{cap:numerica}
como o critério para julgar integradores orbitais.
\end{caixa}

\section{A equação de Kepler}
\label{sec:kepler_eq}

Para avaliar a posição orbital em um instante arbitrário — e não apenas
propagá-la passo a passo — resolve-se a equação de Kepler
\begin{equation}
M = E - e\sin E, \qquad M=\frac{2\pi}{T}(t-t_p),
\label{eq:kepler}
\end{equation}
relacionando a anomalia média $M$ à anomalia excêntrica $E$; a posição no plano
orbital sai de $x=a(\cos E-e)$, $y=a\sqrt{1-e^{2}}\sin E$.

A Eq.~\eqref{eq:kepler} é transcendental. O método de Newton--Raphson,
\begin{equation}
E_{k+1}=E_k-\frac{E_k-e\sin E_k-M}{1-e\cos E_k},
\end{equation}
com o chute $E_0=M+e\sin M$, converge em cerca de quatro iterações para
$e\lesssim 0.1$ e atinge precisão de máquina. A convergência é quadrática
porque $1-e\cos E>0$ para $e<1$, de modo que a derivada nunca se anula.

É esta a efeméride adotada, e a razão da escolha é discutida em
\S\ref{sec:kepleranalitico}: o integrador da
resposta oceânica precisa avaliar a posição da Lua em $t+\Delta t/2$, o que um
propagador passo a passo não oferece de graça.

\subsection{Precessões do plano orbital}

Dois movimentos seculares lentos da órbita lunar são incluídos porque governam
efeitos que a análise harmônica precisa (Capítulo~\ref{cap:harmonica}):

\begin{itemize}
\item o \textbf{nó ascendente} regride com período de \SI{18.613}{\anos}. A
  inclinação da órbita lunar em relação ao \emph{equador} oscila entre
  $\ang{23.44}-\ang{5.15}=\ang{18.3}$ e $\ang{23.44}+\ang{5.15}=\ang{28.6}$;
\item o \textbf{perigeu} avança com período de \SI{8.847}{\anos}, e é ele que
  separa as constituintes N\textsubscript{2} e L\textsubscript{2} de
  M\textsubscript{2}.
\end{itemize}

\begin{alerta}[Um erro real, encontrado por medida]
Manter o nó \emph{fixo} é uma simplificação tentadora, e numa implementação
inicial deste modelo ela deixou a órbita travada por acaso exatamente no máximo
do ciclo nodal. A consequência foi uma superestimação de 12--18\,\% nas
constituintes diurnas — detectada porque a amplitude medida de
O\textsubscript{1} (\num{1.179} em unidades relativas) coincidia com o fator
nodal máximo tabelado ($f_{\max}=\num{1.181}$), e não com a média. Precessões
lentas demais para importar em um mês de simulação importam, sim, para
\emph{calibrar} uma amplitude.
\end{alerta}

\section{Expansão multipolar: de onde vem a maré}
\label{sec:multipolo}

Este é o resultado central da Parte~\ref{parte:fundamentos}. Considere o
potencial gravitacional, por unidade de massa, que um corpo de massa $M$
localizado em $\vect{d}$ (medido do centro da Terra) produz num ponto
$\vect{r}$ do interior ou da superfície terrestre:
\begin{equation}
\Phi(\vect{r})=-\frac{GM}{|\vect{d}-\vect{r}|}.
\end{equation}
Com $\psi$ o ângulo entre $\vect{r}$ e $\vect{d}$, e $r/D\ll 1$, a função
geratriz dos polinômios de Legendre dá
\begin{equation}
\boxed{\;
\Phi(\vect{r})=-\frac{GM}{D}\sum_{n=0}^{\infty}\left(\frac{r}{D}\right)^{n}
P_n(\cos\psi).\;}
\label{eq:multipolo}
\end{equation}

Cada termo tem um significado físico distinto, e a leitura desta série
\emph{é} a teoria da maré:

\begin{description}[leftmargin=1.6cm,style=nextline]
\item[$n=0$: $\;-GM/D$] Constante no espaço. Gradiente nulo: nenhuma força.
  Irrelevante.
\item[$n=1$: $\;-\dfrac{GM r}{D^{2}}\cos\psi$] Gera uma força
  \emph{uniforme} $GM/D^{2}$ dirigida ao corpo — exatamente a aceleração de
  queda livre do centro da Terra. Este termo é \textbf{integralmente cancelado}
  pelo termo de arrasto $-m\ddot{\vect{R}}_0$ da Eq.~\eqref{eq:naoinercial}.
  É o cancelamento que define o problema.
\item[$n=2$: $\;-\dfrac{GM r^{2}}{D^{3}}P_2(\cos\psi)$] O primeiro termo
  \textbf{não} cancelado. É o \emph{potencial de maré}. Todo o resto do
  documento vive aqui.
\item[$n\ge 3$] Menores por fatores $r/D\approx\num{0.0166}$ (Lua) e
  $\num{4.3e-5}$ (Sol). O termo $n=3$ da Lua vale $1.7\,\%$ do $n=2$ e é
  responsável pelas pequenas constituintes de terceiro grau (M\textsubscript{3});
  não é incluído no modelo aqui descrito.
\end{description}

\begin{caixa}[A lição em uma linha]
A maré não é a atração da Lua ($n=1$, cancelada pela queda livre) nem o valor
absoluto do potencial dela no centro da Terra ($n=0$, constante no espaço,
portanto sem gradiente e sem força). É o \textbf{termo quadrupolar} da expansão do
campo do perturbador em torno do centro da Terra — o primeiro que sobrevive à
mudança para o referencial em queda livre. Daí a dependência $D^{-3}$, e daí a
Lua vencer o Sol apesar de atrair a Terra 178 vezes menos do que ele
(\S\ref{sec:solvslua}).
\end{caixa}

\section{Estrutura lagrangiana, hamiltoniana e simplética}
\label{sec:simpletica}

O texto-base é newtoniano, mas dois resultados da formulação hamiltoniana são
usados no Capítulo~\ref{cap:numerica} e convém enunciá-los aqui.

Para um sistema \emph{separável} $H(\vect{q},\vect{p})=T(\vect{p})+V(\vect{q})$
— o caso das órbitas —, o fluxo exato $\varphi_t$ preserva a 2-forma
$\dd\vect{p}\wedge\dd\vect{q}$, isto é, é uma transformação canônica. Um
integrador numérico $\Psi_{\Delta t}$ é \textbf{simplético} quando também
preserva essa forma exatamente, e não apenas até $\ordem{\Delta t^{p}}$.

O que isso compra é dado pela \emph{análise de erro retrógrado}: existe um
hamiltoniano modificado
\begin{equation}
\widetilde{H}=H+\Delta t^{p}H_{p}+\Delta t^{p+1}H_{p+1}+\cdots
\label{eq:hmodificado}
\end{equation}
tal que a trajetória numérica é a trajetória \emph{exata} de $\widetilde{H}$,
a menos de termos exponencialmente pequenos em $1/\Delta t$. Como
$\widetilde{H}$ é conservado exatamente ao longo da trajetória numérica e dista
$\ordem{\Delta t^{p}}$ de $H$, o erro de energia fica \textbf{limitado} por
tempos exponencialmente longos, em vez de crescer secularmente
\cite{hairer2006}.

Métodos não simpléticos — Runge--Kutta explícito, inclusive o RK4 — não
possuem $\widetilde{H}$: o erro de energia tem componente secular. A
Tabela~\ref{tab:integradores} mede exatamente isso, e a
Eq.~\eqref{eq:amplificacao3x} traduz o resultado em erro de maré.

\begin{alerta}[E a armadilha simétrica]
Simplético \textbf{não} implica mais preciso. Velocity-Verlet é simplético e
ainda assim fica ordens de grandeza pior que o RK4 nas escalas de tempo deste
problema, porque é de segunda ordem. A vantagem simplética está no
\emph{formato} do erro, não na sua magnitude.

E há um segundo limite, mais fundamental: a Eq.~\eqref{eq:hmodificado} pressupõe
que o sistema \textbf{seja hamiltoniano}. Um amortecimento $2\gamma\dot h$ o
torna dissipativo, e então não existe $H$ algum a ser modificado: a vantagem
simplética simplesmente deixa de existir. Some-se a isso que uma força
dependente da velocidade quebra a separabilidade $H=T(\vect{p})+V(\vect{q})$ de
que os métodos de \emph{splitting} explícito (Verlet, Yoshida) dependem para
serem explícitos. É justamente o caso da resposta oceânica do
Capítulo~\ref{cap:oscilador}, onde o RK4 volta a ser a escolha correta.
\end{alerta}

% =====================================================================
\chapter{O oscilador forçado amortecido}
\label{cap:oscilador}

O texto-base encerra a seção de marés observando que as marés costeiras reais
superam de longe os \SI{0.54}{\meter} da teoria de equilíbrio e que
``ressonâncias podem afetar a oscilação natural dos corpos d'água''. O modelo
mínimo capaz de formalizar essa observação é o oscilador linear forçado, e ele
merece um capítulo próprio porque toda a Parte~\ref{parte:melhorias} pode ser
lida como a substituição deste modelo de \emph{um} grau de liberdade por um de
$10^{5}$.

\section{Equação, admitância e fator de qualidade}

Uma bacia oceânica com período natural $T_0=2\pi/\omega_0$ forçada pela maré de
equilíbrio $h_{\text{eq}}(t)$ obedece
\begin{equation}
\ddot{h}+2\gamma\dot{h}+\omega_0^{2}h=\omega_0^{2}h_{\text{eq}}(t),
\qquad Q\equiv\frac{\omega_0}{2\gamma}.
\label{eq:oscilador}
\end{equation}
A normalização $\omega_0^2$ do lado direito é escolhida para que a resposta
estática seja $h\to h_{\text{eq}}$ quando $\omega\to 0$.

Para forçamento monocromático $h_{\text{eq}}=H\cos\omega t$, a solução de
regime permanente é $h=A(\omega)H\cos(\omega t-\delta(\omega))$ com
\begin{align}
A(\omega)&=\frac{\omega_0^{2}}
  {\sqrt{(\omega_0^{2}-\omega^{2})^{2}+4\gamma^{2}\omega^{2}}},
\label{eq:admitancia}\\[4pt]
\delta(\omega)&=\atantwo\!\left(2\gamma\omega,\;\omega_0^{2}-\omega^{2}\right)
\;\in\;(0,\pi).
\label{eq:atraso}
\end{align}

Os três regimes de \eqref{eq:admitancia}--\eqref{eq:atraso} descrevem três
tipos de costa reais:

\begin{center}
\begin{tabular}{@{}llll@{}}
\toprule
regime & $A$ & $\delta$ & interpretação física \\
\midrule
$\omega\ll\omega_0$ & $\to 1$ & $\to 0$ & bacia ``rígida'': segue a maré de equilíbrio \\
$\omega\approx\omega_0$ & $\to Q$ & $\pi/2$ & ressonância: amplificação enorme \\
$\omega\gg\omega_0$ & $\to 0$ & $\to\pi$ & bacia lenta: responde \emph{abaixo} do equilíbrio \\
\bottomrule
\end{tabular}
\end{center}

O terceiro regime não é curiosidade formal — há bacias que respondem abaixo do
equilíbrio. Convém, porém, não usar o Taiti como ilustração, apesar de comum: a
maré de Papeete é excepcionalmente pequena ($\sim\SI{0.3}{\meter}$, com
M\textsubscript{2} de poucos centímetros e tábua dominada por
S\textsubscript{2}) por estar praticamente sobre um anfidromo de
M\textsubscript{2} do Pacífico Sul — fenômeno que só o
Capítulo~\ref{cap:anfidromos} explica, e que nada tem a ver com o ramo
$\omega\gg\omega_0$.

\begin{figure}[htbp]
\centering
\begin{tikzpicture}
\begin{axis}[
  width=0.47\textwidth, height=5.2cm,
  xlabel={$\omega/\omega_0$}, ylabel={$A(\omega)$},
  xmin=0, xmax=2.2, ymin=0, ymax=11,
  grid=both, grid style={line width=.2pt,draw=gray!25},
  legend style={font=\scriptsize,draw=none,fill=none,at={(0.98,0.98)},anchor=north east},
  tick label style={font=\scriptsize}, label style={font=\small},
]
\addplot[azulforte,very thick,domain=0.01:2.2,samples=300]
  {1/sqrt((1-x^2)^2 + x^2/(2^2))};
\addlegendentry{$Q=2$}
\addplot[laranja,very thick,domain=0.01:2.2,samples=300]
  {1/sqrt((1-x^2)^2 + x^2/(5^2))};
\addlegendentry{$Q=5$}
\addplot[verde,very thick,domain=0.01:2.2,samples=400]
  {1/sqrt((1-x^2)^2 + x^2/(10^2))};
\addlegendentry{$Q=10$}
\end{axis}
\end{tikzpicture}\hfill
\begin{tikzpicture}
\begin{axis}[
  width=0.47\textwidth, height=5.2cm,
  xlabel={$\omega/\omega_0$}, ylabel={$\delta$ [graus]},
  xmin=0, xmax=2.2, ymin=0, ymax=180,
  ytick={0,45,90,135,180},
  grid=both, grid style={line width=.2pt,draw=gray!25},
  legend style={font=\scriptsize,draw=none,fill=none,at={(0.98,0.30)},anchor=east},
  tick label style={font=\scriptsize}, label style={font=\small},
]
\addplot[azulforte,very thick,domain=0.01:2.2,samples=300]
  {atan2(x/2, 1-x^2)};
\addlegendentry{$Q=2$}
\addplot[laranja,very thick,domain=0.01:2.2,samples=400]
  {atan2(x/5, 1-x^2)};
\addlegendentry{$Q=5$}
\addplot[verde,very thick,domain=0.01:2.2,samples=600]
  {atan2(x/10, 1-x^2)};
\addlegendentry{$Q=10$}
\end{axis}
\end{tikzpicture}
\caption{Admitância \eqref{eq:admitancia} e atraso de fase \eqref{eq:atraso} do
oscilador forçado. À ressonância a amplificação vale $Q$ e o atraso é
exatamente \ang{90}, independentemente de $Q$ — o cruzamento em $\pi/2$ é a
assinatura observável da ressonância, e é mais robusto que o pico de amplitude.}
\label{fig:admitancia}
\end{figure}

\section{Do atraso de fase ao deslocamento espacial do bojo}
\label{sec:deslocamento}

O texto-base observa que a maré alta não ocorre exatamente sob o eixo
Terra--Lua. A tradução quantitativa do atraso temporal em deslocamento angular
requer atenção a um fator 2 que se erra com facilidade.

O padrão de maré é de \emph{grau 2}: a cada rotação relativa da Terra em
relação à Lua, um observador atravessa \textbf{dois} ciclos completos de fase.
Portanto, um atraso de fase $\delta$ no sinal temporal corresponde a um
deslocamento \emph{espacial} do bojo de
\begin{equation}
\Delta\lambda=\frac{\delta}{2}.
\label{eq:deslocamento}
\end{equation}
Uma bacia em ressonância, com $\delta=\ang{90}$, tem o bojo adiantado
\ang{45} em longitude — não \ang{90}.

\section{Transiente, tempo de decaimento e \emph{spin-up}}
\label{sec:spinup}

A solução homogênea de \eqref{eq:oscilador} decai como
$e^{-\gamma t}$, ou seja, com constante de tempo
\begin{equation}
\tau=\frac{1}{\gamma}=\frac{2Q}{\omega_0}=\frac{Q}{\pi}T_0 .
\label{eq:decaimento}
\end{equation}

\begin{alerta}[Uma regra de bolso que custa a figura]
É comum estimar o transiente por ``$Q$ períodos''. A constante de tempo correta
é \eqref{eq:decaimento}, e o critério de $6\tau$ corresponde a
$6Q T_0/\pi\approx\num{1.9}\,Q\,T_0$ — quase o dobro da regra de bolso. Foi
exatamente esse o erro cometido numa primeira implementação, e o sintoma é
característico: a amplitude \emph{cresce} durante os primeiros \num{15} dias
de saída, mascarando por completo o envelope de sizígia/quadratura, que é o
resultado que se queria mostrar. O diagnóstico só ficou limpo quando a constante
de tempo foi derivada em vez de estimada.

O custo precisa ser orçado: bacias com $Q=20$ e $T_0\approx\SI{12.5}{\hour}$
têm $\tau=\SI{3.3}{\day}$, de modo que $6\tau$ são \SI{20}{\day} de
integração descartada.
\end{alerta}

O critério adotado é integrar $6\tau$ antes de $t=0$, deixando
$e^{-6}\approx\num{0.25}\,\%$ de transiente residual — muito abaixo da
incerteza do próprio modelo.

\section{Justificativa física dos parâmetros concentrados}

$T_0$ e $Q$ não são livres: para uma baía retangular de comprimento $L$ e
profundidade $H$, aberta numa extremidade, o modo fundamental de quarto de onda
tem
\begin{equation}
T_0=\frac{4L}{\sqrt{gH}}.
\label{eq:quartodeonda}
\end{equation}
Para a Baía de Fundy, $L\approx\SI{270}{\kilo\meter}$ e
$H\approx\SI{60}{\meter}$ dão $T_0\approx\SI{12.4}{\hour}$, o que já indica por
que a baía é candidata a ressonância. O cálculo cuidadoso, porém, não é o da
baía isolada: Garrett \cite{garrett1972} mostrou que o ressoador é o sistema
acoplado Golfo do Maine--Baía de Fundy, cujo modo livre tem período de
$\approx\SI{13.3}{\hour}$ com $Q\approx 5$ — \emph{próximo}, e não idêntico, ao
período de M\textsubscript{2}. É essa quase-ressonância, e não uma coincidência
exata, a explicação clássica das maiores marés do mundo.

O fator $Q$ mede $2\pi$ vezes a razão entre a energia armazenada e a energia
dissipada \emph{por ciclo}, por radiação para o oceano aberto e por atrito de
fundo. É a grandeza que este modelo de um grau de liberdade \emph{não} consegue
prever a partir da geometria — ela precisa ser ajustada.

Essa é a limitação estrutural, e não acidental, do
Capítulo~\ref{cap:oscilador}; o Capítulo~\ref{cap:lte} a remove.

% =====================================================================
\chapter{Dinâmica de fluidos geofísicos: o mínimo necessário}
\label{cap:fluidos}

Este capítulo estabelece as equações que a Parte~\ref{parte:melhorias} resolve.
A derivação é padrão \cite{gill1982,pedlosky1987}, mas as hipóteses precisam
estar explícitas para que os limites do modelo sejam auditáveis.

\section{Equações de Euler no referencial girante}

Aplicando a Eq.~\eqref{eq:naoinercial} a um elemento de fluido e usando
$\vect{F}=-\nabla p/\rho + \vect{g}_{\text{ef}}$:
\begin{equation}
\pd{\vect{u}}{t}+(\vect{u}\cdot\nabla)\vect{u}
+2\vect{\Omega}\times\vect{u}
=-\frac{1}{\rho}\nabla p-\nabla\Phi+\vect{\mathcal{F}},
\label{eq:euler}
\end{equation}
onde $\Phi$ agrega o geopotencial (gravidade $+$ centrífuga) e o potencial de
maré, e $\vect{\mathcal{F}}$ reúne as forças de atrito.

\section{A aproximação de águas rasas}
\label{sec:aguasrasas}

A onda de maré tem comprimento $\lambda\sim c\,T\sim\SI{9000}{\kilo\meter}$
contra profundidades $H\sim\SI{4}{\kilo\meter}$: a razão de aspecto é
$H/\lambda\sim\num{4e-4}$. Nesse limite a componente vertical de
\eqref{eq:euler} degenera no \textbf{equilíbrio hidrostático}
\begin{equation}
\pd{p}{z}=-\rho g \;\;\Longrightarrow\;\; p(z)=p_a+\rho g(\zeta-z),
\end{equation}
e, com $\rho$ constante (oceano \emph{barotrópico}), o gradiente horizontal de
pressão passa a ser independente de $z$:
$\nabla_h p/\rho=g\nabla_h\zeta$. Consequentemente $\vect{u}_h$ também é
independente de $z$, e integrar a continuidade
$\nabla\cdot\vect{u}=0$ de $z=-H$ a $z=\zeta$ com as condições cinemáticas de
fundo e superfície dá (detalhes no Apêndice~\ref{ap:sw})
\begin{equation}
\pd{\zeta}{t}+\nabla_h\cdot\big[(H+\zeta)\vect{u}\big]=0.
\end{equation}
Como $\zeta/H\sim\num{2.5e-4}$ em oceano aberto, lineariza-se $H+\zeta\to H$.

\begin{caixa}[Barotrópico: o que essa palavra custa]
Supor $\rho$ constante elimina de vez a \textbf{maré interna} — as ondas que se
propagam ao longo da picnoclina, com velocidades de $\sim\SI{3}{\meter\per
\second}$ em vez de \SI{200}{\meter\per\second}. Essas ondas levam consigo
$\sim\SI{1}{\tera\watt}$, algo entre 25 e 30\,\% da dissipação total. Um modelo
barotrópico não pode resolvê-las; pode apenas \emph{parametrizar} a perda de
energia que elas representam, e é exatamente o que faz o termo $r\vect{u}$ do
Capítulo~\ref{cap:dissipacao}.
\end{caixa}

\section{Ondas em fluido girante}
\label{sec:ondas}

Linearizando as equações de águas rasas com $f=2\Omega\sin\phi$ constante
(plano-$f$) obtém-se a relação de dispersão das \textbf{ondas de Poincaré}
\begin{equation}
\sigma^{2}=f^{2}+gH\,(k^{2}+l^{2}),
\label{eq:poincare}
\end{equation}
que mostra de imediato um fato decisivo: ondas com $\sigma<f$ \textbf{não se
propagam}. A latitude crítica de M\textsubscript{2}, onde $f=\sigma_{M2}$, é
$\phi\approx\ang{74.5}$: acima dela não há ondas de Poincaré livres nessa
frequência, e a maré semidiurna deixa de se propagar como onda de
gravidade-inércia de oceano aberto. Isso \textbf{não} quer dizer que não haja
maré semidiurna ali — ondas de Kelvin costeiras e ondas topográficas de
plataforma não têm latitude crítica, e são elas que levam M\textsubscript{2} ao
Ártico.

Existe, além disso, uma solução que \eqref{eq:poincare} não captura: com uma
costa em $y=0$ e impondo $v\equiv 0$, resta a \textbf{onda de Kelvin}
\begin{equation}
\zeta=\zeta_0\,e^{-y/L_R}\cos(kx-\sigma t),
\qquad c=\sqrt{gH},\qquad
L_R=\frac{\sqrt{gH}}{f},
\label{eq:kelvin}
\end{equation}
não dispersiva, presa à costa, decaindo na escala do \textbf{raio de deformação
de Rossby} $L_R$ e propagando-se sempre com a costa \emph{à direita} no
hemisfério norte. Em \ang{45} de latitude,
$L_R=\SI{198}{\meter\per\second}/\SI{1.03e-4}{\per\second}\approx
\SI{1900}{\kilo\meter}$.

A onda de Kelvin é o objeto central do Capítulo~\ref{cap:anfidromos}: são duas
delas, incidente e refletida, que geram um ponto anfidrômico.

\section{Análise de escala: Coriolis não é perturbação}
\label{sec:escala}

Antes de resolver qualquer coisa, convém saber quais termos importam. Para a
maré M\textsubscript{2} em \ang{45} de latitude:
\begin{equation}
\frac{f}{\sigma_{M2}}
=\frac{\SI{1.031e-4}{\per\second}}{\SI{1.405e-4}{\per\second}}
=\num{0.73}.
\end{equation}

O termo de Coriolis é da mesma ordem que a aceleração local: é um efeito de
\textbf{primeira ordem}, não uma correção. Já a advecção
$(\vect{u}\cdot\nabla)\vect{u}$ tem número de Rossby
\begin{equation}
\mathrm{Ro}=\frac{U}{fL}\sim
\frac{\SI{0.05}{\meter\per\second}}
{\SI{1e-4}{\per\second}\times\SI{1e6}{\meter}}
=\num{5e-4},
\end{equation}
o que justifica a linearização em oceano aberto. Em estuários e plataformas,
com $U\sim\SI{1}{\meter\per\second}$ e $L\sim\SI{10}{\kilo\meter}$,
$\mathrm{Ro}\sim 1$: é ali que nascem as constituintes de águas rasas
(M\textsubscript{4}, M\textsubscript{6}), e é ali que a linearização deixa de
valer.

\begin{observacao}
Note a coerência interna do argumento: o modelo aqui descrito lineariza a
advecção \emph{e} não resolve estuários. Os dois limites são o mesmo limite. Um
modelo que resolvesse estuários numa grade fina precisaria manter o termo
advectivo.
\end{observacao}

% =====================================================================
\part{A força de maré e a teoria de equilíbrio}
\label{parte:mare}
% =====================================================================

\chapter{A força de maré}
\label{cap:forcamare}

\section{Derivação no referencial geocêntrico}

Seguindo o texto-base (Eq.~5.48--5.51), considere uma parcela d'água de massa
$m$ na superfície da Terra, sob a ação da gravidade terrestre e da atração de
um corpo perturbador de massa $M$ (a Lua, para fixar ideias). No referencial
inercial,
\begin{equation}
m\,\ddot{\vect{r}}\,' _m=-\frac{GmM_\oplus}{r^{2}}\uvec{e}_r
-\frac{GmM}{R^{2}}\uvec{e}_R ,
\tag{T\&M~5.48}
\label{eq:548}
\end{equation}
enquanto o centro da Terra, tratado como partícula livre no campo do corpo,
obedece
\begin{equation}
M_\oplus\,\ddot{\vect{r}}\,'_\oplus=-\frac{GM_\oplus M}{D^{2}}\uvec{e}_D .
\tag{T\&M~5.49}
\label{eq:549}
\end{equation}

Subtraindo \eqref{eq:549} de \eqref{eq:548} — que é precisamente aplicar o
termo de arrasto $-m\ddot{\vect{R}}_0$ da Eq.~\eqref{eq:naoinercial} —
obtém-se a aceleração medida no referencial geocêntrico:
\begin{equation}
\ddot{\vect{r}}=
\underbrace{-\frac{GM_\oplus}{r^{2}}\uvec{e}_r}_{\text{gravidade local}}
\;\underbrace{-\;GM\left(\frac{\uvec{e}_R}{R^{2}}
-\frac{\uvec{e}_D}{D^{2}}\right)}_{\text{força de maré}} .
\tag{T\&M~5.50--5.51}
\label{eq:551}
\end{equation}

A força de maré é, portanto, uma \textbf{diferença de atrações}, não uma
atração: a Lua puxa a água próxima mais forte do que puxa o centro da Terra, e
o centro da Terra mais forte do que a água do lado oposto. Os dois bojos
antipodais saem daí, e não de nenhuma ``força centrífuga do sistema
Terra--Lua'' — explicação frequente e incorreta, já que o resultado vale
igualmente para um corpo em queda livre radial, sem rotação alguma.

\begin{alerta}[A convenção de sinais que mais custa tempo]
Pela figura do texto-base (Fig.~5-10a):
\begin{center}
\begin{tabular}{@{}ll@{}}
\toprule
símbolo & aponta de $\to$ para \\
\midrule
$\uvec{e}_r$ & centro da Terra $\to$ parcela \\
$\uvec{e}_R$ & \textbf{Lua} $\to$ parcela \\
$\uvec{e}_D$ & \textbf{Lua} $\to$ centro da Terra \\
\bottomrule
\end{tabular}
\end{center}
Inverter $\uvec{e}_D$ transforma a \emph{diferença} de \eqref{eq:551} numa
\emph{soma}: no ponto sublunar o resultado passa a valer $\approx 2GM/D^{2}$ em
vez de $2GMr/D^{3}$, ou seja, é maior por um fator $D/r\approx 60$, e troca de
sinal no lado antipodal. O erro não se anuncia — a expressão continua com
aparência correta — e só aparece quando se compara a forma exata com a expandida
\eqref{eq:expandida}, que é imune a ele: ali $\uvec{d}$ entra
\emph{quadraticamente}, e inverter o versor não muda nada.
\end{alerta}

Para eliminar a ambiguidade, convém escrever \eqref{eq:551} na forma
equivalente e autoexplicativa
\begin{equation}
\vect{a}_T=\underbrace{\frac{GM(\vect{d}-\vect{r})}{|\vect{d}-\vect{r}|^{3}}}
_{\text{atração no ponto}}
-\underbrace{\frac{GM\vect{d}}{|\vect{d}|^{3}}}_{\text{atração no centro}},
\label{eq:exata}
\end{equation}
com $\vect{d}$ o vetor Terra~$\to$~corpo, que é a convenção natural das
efemérides. A identidade com \eqref{eq:551} decorre de
$\uvec{e}_R=(\vect{r}-\vect{d})/|\vect{r}-\vect{d}|$ e
$\uvec{e}_D=-\vect{d}/|\vect{d}|$.

\begin{figure}[htbp]
\centering
\begin{tikzpicture}[scale=1.0,>=Latex]
  % Terra
  \draw[fill=azulclaro,draw=azulforte,thick] (0,0) circle (1.35);
  \fill[azulforte] (0,0) circle (1.5pt);
  \node[below left] at (0,0) {\small $\oplus$};
  % parcela
  \coordinate (P) at (35:1.35);
  \fill[laranja] (P) circle (2pt);
  \node[above right=1pt] at (P) {\small $m$};
  % Lua
  \coordinate (L) at (7.6,0);
  \draw[fill=gray!35,draw=gray!60] (L) circle (0.32);
  \node[below=6pt] at (L) {\small Lua, $M$};
  % vetores
  \draw[->,thick,azulforte] (0,0) -- (P) node[midway,above left=-2pt] {\small $\vect{r}$};
  \draw[->,thick,laranja] (L) -- (P) node[midway,above=1pt] {\small $R\,\uvec{e}_R$};
  \draw[->,thick,verde] (L) -- (0,0) node[midway,below=2pt] {\small $D\,\uvec{e}_D$};
  \draw[dashed,gray] (0,0) -- (8.2,0);
  \draw (1.1,0) arc (0:35:1.1);
  \node at (0.95,0.28) {\small $\theta$};
\end{tikzpicture}
\caption{Geometria da Eq.~\eqref{eq:551}. Os versores $\uvec{e}_R$ e
$\uvec{e}_D$ apontam \emph{a partir da Lua}; $\uvec{e}_r$ e $\theta$ são
medidos no centro da Terra. A figura não está em escala:
$D/r_\oplus\approx 60$.}
\label{fig:geometria}
\end{figure}

\section{Expansão em $r/D$ e o tensor de maré}

Com $r/D=\num{0.0166}$ para a Lua, expandir \eqref{eq:exata} e reter o primeiro
termo não nulo dá, no ponto sublunar e em quadratura,
\begin{align}
F_{Tx}&=+\frac{2GmM r}{D^{3}}, \tag{T\&M~5.52}\label{eq:552}\\
F_{Ty}&=-\frac{GmM r}{D^{3}}, \tag{T\&M~5.53}
\end{align}
e, para um ponto qualquer a um ângulo $\theta$ do eixo Terra--Lua,
\begin{equation}
F_{Tx}=\frac{2GmMr\cos\theta}{D^{3}},\qquad
F_{Ty}=-\frac{GmMr\sin\theta}{D^{3}} .
\tag{T\&M~5.54a,b}
\label{eq:554}
\end{equation}

A forma vetorial tridimensional — que é a implementada — é
\begin{equation}
\boxed{\;\vect{a}_T(\vect{r})=\frac{GM}{D^{3}}
\left[3(\vect{r}\cdot\uvec{d})\,\uvec{d}-\vect{r}\right].\;}
\label{eq:expandida}
\end{equation}
Com $\uvec{d}=\uvec{x}$ e $\vect{r}=(x,y,z)$, ela devolve
$(GM/D^{3})(2x,-y,-z)$, isto é, \eqref{eq:552} e~(5.53) simultaneamente.

\begin{observacao}[Leitura tensorial]
A Eq.~\eqref{eq:expandida} é $a_{T,i}=T_{ij}r_j$ com o \emph{tensor de maré}
\begin{equation}
T_{ij}=-\pd{^{2}\Phi}{x_i\partial x_j}\bigg|_{\vect{r}=0}
=\frac{GM}{D^{3}}\left(3\hat d_i\hat d_j-\delta_{ij}\right),
\end{equation}
cujo traço é $\;\mathrm{tr}\,T=\frac{GM}{D^{3}}(3-3)=0$. O traço nulo é a
equação de Laplace $\nabla^{2}\Phi=0$: não há massa no ponto considerado. É por
isso que a compressão em duas direções vem sempre acompanhada de \emph{uma}
distensão. O traço nulo obriga apenas que os três autovalores somem zero;
combinado com a simetria axial em torno de $\uvec{d}$, isso fixa a razão
$(+2,-1,-1)$ — obrigatória, e não uma escolha do modelo.
\end{observacao}

\begin{figure}[htbp]
\centering
\begin{tikzpicture}[scale=1.5,>=Latex]
  \draw[fill=azulclaro,draw=azulforte,thick] (0,0) circle (1);
  \foreach \t in {0,15,...,345}{
    \pgfmathsetmacro{\ax}{0.42*2*cos(\t)}
    \pgfmathsetmacro{\ay}{-0.42*sin(\t)}
    \draw[->,azulforte,thick]
      ({cos(\t)},{sin(\t)}) -- ++({\ax},{\ay});
  }
  \draw[dashed,gray] (-2.2,0) -- (2.9,0);
  \node[gray] at (2.65,0.22) {\small $\to$ Lua};
  \node at (0,-2.15) {};
\end{tikzpicture}
\caption{Campo da Eq.~\eqref{eq:554} na superfície: distensão de magnitude
$2GMr/D^{3}$ ao longo do eixo Terra--Lua, compressão de magnitude $GMr/D^{3}$
perpendicularmente a ele. As componentes tangenciais — e não as radiais — são
as que efetivamente movem a água, porque a componente radial compete com
$g$ e vale $\sim 10^{-7}g$.}
\label{fig:campo}
\end{figure}

\section{O erro da expansão, medido}

A diferença relativa entre \eqref{eq:exata} e \eqref{eq:expandida}, calculada
numericamente sobre a superfície terrestre para a Lua, fica entre
\SI{1.7}{\percent} e \SI{2.5}{\percent} conforme o ângulo — exatamente
$\ordem{r/D}$, como o texto-base afirma. Dobrar $r/D$ dobra o erro, o que
confirma a ordem, e não apenas a magnitude.

\begin{caixa}[Metodologia: por que manter as duas formas]
A forma exata \eqref{eq:exata} é tão barata quanto a expandida e mais precisa;
manter as duas implementadas não serve ao cálculo, serve à \emph{auditoria}. É a comparação
entre elas que transforma ``a expansão é $\ordem{r/D}$'' de afirmação de
livro-texto em resultado medido — e foi ela que revelou o erro de sinal em
$\uvec{e}_D$. Uma implementação que contenha apenas a versão que de fato usa não tem
como detectar esse tipo de erro.
\end{caixa}

% =====================================================================
\chapter{O potencial de maré e a maré de equilíbrio}
\label{cap:potencial}

\section{Do campo de forças ao potencial}

O texto-base trabalha só com forças. Para obter a \emph{forma da superfície} é
preciso o potencial, que se obtém integrando $\vect{F}=-\nabla U$ a partir de
\eqref{eq:554} — ou, equivalentemente, lendo o termo $n=2$ da expansão
multipolar \eqref{eq:multipolo}:
\begin{equation}
\boxed{\;
\frac{U_T}{m}=-\frac{GM}{2D^{3}}\left[3(\vect{r}\cdot\uvec{d})^{2}-r^{2}\right]
=-\frac{GMr^{2}}{2D^{3}}\left(3\cos^{2}\psi-1\right)
=-\frac{GMr^{2}}{D^{3}}P_2(\cos\psi),\;}
\label{eq:potencial}
\end{equation}
com $\psi$ o ângulo entre a parcela e o eixo Terra--corpo. Derivando de volta
recupera-se $-\partial U/\partial x=2GMx/D^{3}$ e
$-\partial U/\partial y=-GMy/D^{3}$, o que fecha a consistência com
\eqref{eq:552}--(T\&M~5.53), o que se verifica também por diferenciação
numérica direta.

\section{A superfície de equilíbrio}

A hipótese de \textbf{equilíbrio} é que a superfície do oceano é uma
equipotencial do campo total. Se assim for, a elevação satisfaz
$g\,h+U_T/m=\text{const}$, ou seja $h=-U_T/(mg)$:
\begin{equation}
h_{\text{eq}}(\psi)=\frac{GMr^{2}}{2gD^{3}}\left(3\cos^{2}\psi-1\right)
=\frac{1}{2}\frac{M}{M_\oplus}\frac{r^{4}}{D^{3}}\left(3\cos^{2}\psi-1\right),
\label{eq:heq}
\end{equation}
onde a segunda forma usou $g=GM_\oplus/r^{2}$ e elimina $G$ e $g$ do cálculo —
detalhe numérico útil, porque $G$ é a constante fundamental menos bem conhecida.

A diferença entre maré alta ($\psi=0$) e baixa ($\psi=\ang{90}$) — a
\emph{excursão} da maré, e não a sua amplitude — reproduz a Eq.~5.55 do
texto-base:
\begin{equation}
h=\frac{3GMr^{2}}{2gD^{3}} .
\tag{T\&M~5.55}
\label{eq:555}
\end{equation}

O fator $(3\cos^{2}\psi-1)$ tem máximos em $\psi=0$ \emph{e} $\psi=\ang{180}$:
os dois bojos antipodais, e portanto \textbf{duas} marés altas por dia lunar.
Este é o resultado contraintuitivo que Galileu não conseguiu explicar e que
mede a distância entre uma teoria correta e uma plausível.

\section{As três espécies: decomposição em harmônicos esféricos}
\label{sec:especies}

A Eq.~\eqref{eq:potencial} está escrita em termos de $\psi$, ângulo que depende
simultaneamente da posição do observador e da posição instantânea do corpo. Para
separar geografia de astronomia — o que é indispensável à análise harmônica —
usa-se
\begin{equation}
\cos\psi=\sin\phi\sin\delta+\cos\phi\cos\delta\cos\mathcal{H},
\end{equation}
com $\phi$ a latitude do observador, $\delta$ a declinação do corpo e $H$ o
ângulo horário — grafado $\mathcal{H}$ para não colidir com a profundidade
$H$ da Parte~\ref{parte:melhorias}. O teorema da adição dos harmônicos esféricos
(Apêndice~\ref{ap:legendre}) então fatora $P_2(\cos\psi)$ em três termos:
\begin{equation}
\boxed{
\begin{aligned}
P_2(\cos\psi)=\;&P_2(\sin\phi)\,P_2(\sin\delta) &&\text{(espécie 0: longo período)}\\
+\;&\tfrac13 P_2^{1}(\sin\phi)\,P_2^{1}(\sin\delta)\cos\mathcal{H} &&\text{(espécie 1: diurna)}\\
+\;&\tfrac1{12}P_2^{2}(\sin\phi)\,P_2^{2}(\sin\delta)\cos 2\mathcal{H} &&\text{(espécie 2: semidiurna)}
\end{aligned}}
\label{eq:especies}
\end{equation}
com $P_2(x)=\tfrac12(3x^{2}-1)$, $P_2^{1}(x)=3x\sqrt{1-x^{2}}$ e
$P_2^{2}(x)=3(1-x^{2})$.

As dependências em latitude que daí resultam são as formas geodésicas padrão
com que cada constituinte é escrita:
\begin{center}
\begin{tabular}{@{}llll@{}}
\toprule
espécie & dependência em $\phi$ & modulação temporal & exemplos \\
\midrule
0 (longo período) & $\tfrac12(3\sin^{2}\phi-1)$ & $\cos(\sigma t)$ & Mf, Mm, Ssa \\
1 (diurna) & $\sin 2\phi$ & $\cos(\sigma t+\lambda)$ & K\textsubscript{1}, O\textsubscript{1}, P\textsubscript{1} \\
2 (semidiurna) & $\cos^{2}\phi$ & $\cos(\sigma t+2\lambda)$ & M\textsubscript{2}, S\textsubscript{2}, N\textsubscript{2} \\
\bottomrule
\end{tabular}
\end{center}

Três consequências físicas imediatas, todas observáveis:

\begin{itemize}
\item as marés \textbf{semidiurnas} anulam-se nos polos ($\cos^{2}\phi=0$) e são
  máximas no equador;
\item as \textbf{diurnas} anulam-se no equador \emph{e} nos polos, com máximo em
  \ang{45} — e, mais importante, seu coeficiente contém
  $P_2^{1}(\sin\delta)\propto\sin 2\delta$, ou seja, \textbf{desaparecem quando
  a declinação é nula}. A desigualdade diurna é um efeito de declinação;
\item as de \textbf{longo período} têm um coeficiente que muda de sinal em
  $\phi=\ang{35.26}$, o que explica a correção de fase de \ang{180} aplicada a
  essa espécie na análise harmônica (\S\ref{sec:correcaofase}).
\end{itemize}

\section{Números reproduzidos}
\label{sec:numeros}

Calculados com as constantes exatas adotadas pelo texto-base, e não com as
constantes modernas, que diferem em $\sim\SI{0.1}{\percent}$ e deslocam $h$ na
terceira casa decimal:

\begin{table}[htbp]
\centering
\caption{Reprodução dos resultados do texto-base. A concordância na terceira
casa decimal valida simultaneamente a álgebra e as constantes.}
\label{tab:numeros}
\begin{tabular}{@{}lrr@{}}
\toprule
grandeza & calculado & texto-base \\
\midrule
$h$ lunar (excursão), Eq.~\eqref{eq:555} & \SI{0.5377}{\meter} & \SI{0.54}{\meter} \\
$h$ solar (excursão) & \SI{0.2461}{\meter} & --- \\
razão Sol/Lua (efeito de maré) & \num{0.4577} & \num{0.46} \\
razão Sol/Lua (atração direta) & \num{178} & $\sim\!\num{175}$ \\
$r/D$ (Lua) & \num{0.0166} & \num{0.02} \\
período de M\textsubscript{2} & \SI{12}{\hour}\,\SI{25.2}{\minute} & \SI{12}{\hour}\,\SI{26}{\minute} \\
erro da expansão & \SIrange{1.7}{2.5}{\percent} & $\ordem{r/D}$ \\
\bottomrule
\end{tabular}
\end{table}

\subsection{Sol contra Lua: o argumento decisivo}
\label{sec:solvslua}

O contraste entre as duas primeiras razões da Tabela~\ref{tab:numeros} é o
coração da seção do livro:
\begin{equation}
\frac{F_{\text{Sol}}}{F_{\text{Lua}}}
=\frac{M_\odot}{M_{\text{Lua}}}\left(\frac{D_{\text{Lua}}}{D_\odot}\right)^{2}
=\num{178},
\qquad
\frac{h_\odot}{h_{\text{Lua}}}
=\frac{M_\odot}{M_{\text{Lua}}}\left(\frac{D_{\text{Lua}}}{D_\odot}\right)^{3}
=\num{0.458}.
\end{equation}

\textbf{O Sol atrai a Terra 178 vezes mais forte que a Lua e ainda assim produz
menos da metade da maré.} A única diferença entre as duas expressões é um fator
$D_{\text{Lua}}/D_\odot=\num{2.6e-3}$ — a marca registrada do quadrupolo. Quem
entendeu esta linha entendeu \S\ref{sec:multipolo}.

\subsection{Sizígia, quadratura e o erratum do texto}

Somando as duas marés em fase e em oposição:
\begin{equation}
h_{\text{sizígia}}=(1+\num{0.4577})\,h_{\text{Lua}}=\num{1.46}\,h,
\qquad
h_{\text{quadratura}}=(1-\num{0.4577})\,h_{\text{Lua}}=\num{0.54}\,h ,
\end{equation}
com período de batimento igual a meio mês sinódico, \SI{14.77}{\day}.

\begin{alerta}[Erratum no texto-base]
A p.~203 afirma: \emph{``The maximum tide, which occurs every 2 weeks, should be
$1.46h=\SI{0.83}{\meter}$ for the spring tides.''} Com
$h=\SI{0.54}{\meter}$, porém, $1.46\times 0.54=\SI{0.79}{\meter}$, não
\SI{0.83}{\meter}. O fator e o valor impressos são mutuamente inconsistentes
($0.83/1.46=\SI{0.57}{\meter}$, que não é o $h$ do próprio exemplo). O valor
aritmeticamente correto é \SI{0.785}{\meter}, e é o adotado aqui.
\end{alerta}

\section{Terra elástica: números de Love}
\label{sec:love}

A Terra sólida também se deforma sob a mesma força. Dois números adimensionais
de grau 2 descrevem o efeito: $h_2\approx\num{0.61}$, a razão entre o
deslocamento radial do solo e a maré de equilíbrio, e $k_2\approx\num{0.30}$, a
razão entre o potencial adicional gerado pela redistribuição de massa e o
potencial forçante.

Um marégrafo está preso ao solo, e mede a diferença entre a superfície do mar e
o solo que também subiu. A maré \emph{medida} é portanto reduzida pelo fator
diminutivo
\begin{equation}
\gamma_2\equiv 1+k_2-h_2\approx\num{0.69}.
\label{eq:love}
\end{equation}
O mesmo fator aparece na Parte~\ref{parte:melhorias} multiplicando a forçante
das Equações de Maré de Laplace. A convenção adotada aqui é deliberada e vale
registrá-la: nos resultados de maré de equilíbrio o fator é \emph{omitido}, para
reproduzir os números clássicos de Terra rígida; no modelo dinâmico ele é
\emph{aplicado}, porque ali o objetivo é reproduzir o oceano.

% =====================================================================
\chapter{Primeiro nível dinâmico: bacias como osciladores}
\label{cap:bacias}

A teoria de equilíbrio prevê no máximo \SI{0.54}{\meter} — e só no equador, com
a Lua no zênite; menos em qualquer outro lugar, pela dependência em latitude de
\eqref{eq:especies}. Marés reais vão de
\SI{0.2}{\meter} no Taiti a \SI{16}{\meter} na Baía de Fundy — duas ordens de
grandeza de espalhamento. O primeiro modelo capaz de produzir esse
espalhamento é o do Capítulo~\ref{cap:oscilador} aplicado estação a estação:
cada bacia é um oscilador com seu $T_0$ e seu $Q$, forçado pela maré de
equilíbrio calculada no local.

\section{Resultados e sua leitura}

\begin{table}[htbp]
\centering
\caption{Amplitude modelada pelo oscilador de um grau de liberdade contra
observação. O modelo acerta a ordem de
grandeza em baixa latitude e subestima por 2--3$\times$ nas bacias
ressonantes de alta latitude.}
\label{tab:estacoes}
\begin{tabular}{@{}lrrl@{}}
\toprule
estação & modelado & observado & caráter da bacia \\
\midrule
Rio de Janeiro & \SI{1.11}{\meter} & \SI{1.20}{\meter} & amplificação modesta \\
Santos & \SI{1.15}{\meter} & \SI{1.40}{\meter} & micro-maré \\
Belém & \SI{2.99}{\meter} & \SI{3.50}{\meter} & foz amplificadora \\
São Luís & \SI{4.08}{\meter} & \SI{6.60}{\meter} & plataforma ampla \\
\textbf{Burntcoat Head (Fundy)} & \textbf{\SI{5.2}{\meter}} & \textbf{\SI{16.0}{\meter}} & quase ressonante \\
\bottomrule
\end{tabular}
\end{table}

O padrão da Tabela~\ref{tab:estacoes} é informativo: o erro cresce
sistematicamente com a latitude e com o grau de ressonância. Duas causas
físicas, ambas identificáveis:

\begin{enumerate}
\item \textbf{A forçante local é fraca em alta latitude.} Pela
  Eq.~\eqref{eq:especies}, a componente semidiurna varia com $\cos^{2}\phi$: a
  \ang{45} ela vale \SI{50}{\percent} da equatorial, e a maré de equilíbrio
  \emph{total} no local cerca de \SI{66}{\percent}. Um ponto em Fundy jamais
  chega perto do ponto sublunar.
\item \textbf{A energia não vem de onde o modelo supõe.} Os \SI{16}{\meter} de
  Fundy não são produzidos pelo forçamento astronômico \emph{local}; são o
  resultado da co-oscilação da baía com o sistema anfidrômico do Atlântico
  Norte inteiro \cite{garrett1972}. A baía é um ressonador excitado pela maré
  oceânica em sua boca, não pelo potencial de maré em seu interior.
\end{enumerate}

\begin{caixa}[Diagnóstico honesto de um modelo]
A subestimação em Fundy \textbf{não é erro de ajuste} — nenhum valor de $Q$ a
corrige sem estragar o resto, porque a via de excitação está ausente da
formulação. É física faltante. Isso muda o que se deve fazer a seguir: não
recalibrar, e sim trocar de modelo. Os parâmetros $T_0$ e $Q$ de cada bacia
são, portanto, parâmetros concentrados ajustados, e é assim que devem ser
declarados — apresentá-los como se fossem deriváveis da geometria falsearia a
natureza do modelo.
\end{caixa}

\section{O que este nível já entrega}

Apesar do limite, o modelo de bacia produz três resultados que a teoria de
equilíbrio pura não produz, e todos são verificáveis:

\begin{itemize}
\item \textbf{amplificação e supressão} — uma bacia lenta ($T_0\gg T_{M2}$)
  responde \emph{abaixo} do equilíbrio, $A<1$, comportamento que só existe no
  regime $\omega\gg\omega_0$ da Fig.~\ref{fig:admitancia} e que a teoria de
  equilíbrio não admite;
\item \textbf{atraso de fase} — via \eqref{eq:deslocamento}, o bojo dinâmico
  aparece deslocado do eixo Terra--Lua, reproduzindo qualitativamente a
  Fig.~5-13 do texto-base;
\item \textbf{envelope de sizígia/quadratura} — o batimento de \SI{14.77}{\day}
  emerge da soma das forçantes lunar e solar, sem ser codificado.
\end{itemize}

% =====================================================================
\chapter{Por que a teoria de equilíbrio tem de falhar}
\label{cap:falha}

O capítulo anterior mostrou que a teoria de equilíbrio erra em \emph{magnitude}.
Este mostra algo mais forte: ela erra em \textbf{forma funcional}, e o
argumento é de uma linha.

\section{A corrida que a onda perde}
\label{sec:corrida}

A maré de equilíbrio pressupõe que o oceano assume instantaneamente a forma da
equipotencial. Isso exige que a informação — isto é, a onda de gravidade —
viaje ao menos tão rápido quanto o padrão forçante se desloca. Compare:
\begin{align}
\text{onda de gravidade em água rasa:}\quad
c&=\sqrt{gH}=\sqrt{\num{9.81}\times\num{4000}}
=\SI{198}{\meter\per\second},\\
\text{ponto sublunar, varrendo o solo:}\quad
v&=\frac{2\pi a}{T_{\text{dia lunar}}}
=\frac{2\pi\times\SI{6.371e6}{\meter}}{\SI{89428}{\second}}
=\SI{448}{\meter\per\second}.
\end{align}

\begin{caixa}[A onda perde por $\num{2.3}\times$]
Não existe bojo hidrostático sob a Lua porque o oceano \textbf{não consegue
construí-lo}. Para acompanhar o ponto sublunar seria preciso
$H=v^{2}/g\approx\SI{20}{\kilo\meter}$ de profundidade — cinco vezes o oceano
real. O que existe não é um bojo atrasado: é uma \emph{onda forçada} muito
abaixo da sua velocidade de fase natural, que reflete nas costas e ressoa nos
modos de cada bacia.
\end{caixa}

\begin{alerta}[Qual velocidade comparar]
A velocidade relevante do ponto sublunar é a \textbf{relativa ao solo},
$2\pi a/T_{\text{dia lunar}}=\SI{448}{\meter\per\second}$, e não
$\Omega a=\SI{465}{\meter\per\second}$ — esta é a velocidade de um ponto do
equador no espaço inercial, que não é com quem a onda compete. A diferença é de
\SI{4}{\percent} e não muda a conclusão, mas a versão errada aparece com
frequência na literatura de divulgação.
\end{alerta}

\section{Coriolis não é uma correção}

Retomando \S\ref{sec:escala}: a \ang{45} de latitude, $f/\sigma_{M2}=\num{0.73}$.
Um termo dessa magnitude não pode ser tratado perturbativamente. E ele
\textbf{muda a topologia da solução}: sob Coriolis a onda forçada organiza-se em
sistemas rotativos em torno de pontos de amplitude nula, como se demonstra no
Capítulo~\ref{cap:anfidromos}.

A teoria de equilíbrio não tem como produzir esses pontos. Nela a maré está
sempre em fase com o potencial, de modo que a fase assume apenas dois valores
(\ang{0} e \ang{180}) e \textbf{nunca circula}. A verificação é direta:
aplicando ao campo de equilíbrio o mesmo detector de anfidromos usado no campo
dinâmico (\S\ref{sec:topologia}), ele não encontra nenhum.

\begin{caixa}[Metodologia: testar a ausência]
Verificar que um modelo \emph{não} produz um fenômeno é tão importante quanto
verificar que o outro produz. Sem esse teste, a detecção de anfidromos no
modelo dinâmico poderia ser um artefato do detector. Com ele, os dois modelos
são submetidos ao mesmo instrumento e respondem de forma qualitativamente
distinta — o que torna a diferença atribuível à física, e não ao diagnóstico.
\end{caixa}

\section{O que falta acrescentar}

O diagnóstico dos Capítulos~\ref{cap:bacias} e~\ref{cap:falha} determina a
lista de melhorias que a Parte~\ref{parte:melhorias} implementa:

\begin{center}
\small
\begin{tabular}{@{}lll@{}}
\toprule
sintoma & causa & correção implementada \\
\midrule
sem anfidromos & Coriolis ausente & LTE com $f=2\Omega\sin\phi$ (Cap.~\ref{cap:lte}) \\
onda instantânea & propagação ausente & águas rasas, $c=\sqrt{gH}$ (Cap.~\ref{cap:lte}) \\
costas invisíveis & geometria ausente & batimetria ETOPO + máscara (\S\ref{sec:batimetria}) \\
amplitude sem limite & dissipação ausente & atrito de fundo + maré interna (Cap.~\ref{cap:dissipacao}) \\
geoide rígido & auto-atração ausente & SAL, $\beta$ (Cap.~\ref{cap:sal}) \\
$Q$ arbitrário & ressonância imposta à mão & modos emergem da bacia real \\
\bottomrule
\end{tabular}
\end{center}

% =====================================================================
\part{Fundamento teórico e formulação das melhorias}
\label{parte:melhorias}
% =====================================================================

\chapter{As Equações de Maré de Laplace}
\label{cap:lte}

Este capítulo formula a primeira e mais importante das melhorias: substituir a
hipótese de equilíbrio pela hidrodinâmica. É o salto de 1687 para 1776 — de
Newton para Laplace.

\section{Derivação}

Parta das equações de águas rasas (\S\ref{sec:aguasrasas}) na esfera girante,
com $a$ o raio da Terra, $\lambda$ a longitude, $\phi$ a latitude,
$\vect{u}=(u,v)$ a velocidade média na vertical (o transporte é $H\vect{u}$) e
$\zeta$ a elevação. A
linearização é justificada por $\mathrm{Ro}\sim\num{5e-4}$ e
$\zeta/H\sim\num{2.5e-4}$ (\S\ref{sec:escala}). Escrevendo os operadores
métricos em coordenadas esféricas:

\begin{empheq}[box=\fbox]{align}
\pd{\zeta}{t}&=-\frac{1}{a\cos\phi}
\left[\pd{(Hu)}{\lambda}+\pd{(Hv\cos\phi)}{\phi}\right],
\label{eq:lte_cont}\\[4pt]
\pd{u}{t}&=\;\;f v-\frac{g}{a\cos\phi}\pd{\zeta'}{\lambda}
-\frac{C_d\,|\vect{u}|\,u}{H}-r\,u,
\label{eq:lte_u}\\[4pt]
\pd{v}{t}&=-f u-\frac{g}{a}\pd{\zeta'}{\phi}
-\frac{C_d\,|\vect{u}|\,v}{H}-r\,v,
\label{eq:lte_v}
\end{empheq}

com $f=2\Omega\sin\phi$ e o \textbf{desvio de equipotencial}
\begin{equation}
\boxed{\;\zeta'=(1-\beta)\,\zeta-\zeta_{\text{eq}}.\;}
\label{eq:desvio}
\end{equation}

A Eq.~\eqref{eq:desvio} concentra duas das melhorias e merece leitura cuidadosa.
O que impulsiona a água não é o gradiente da superfície do mar, e sim o
gradiente da superfície \emph{relativa à equipotencial instantânea}: um oceano
cuja superfície já coincidisse com a equipotencial não sentiria força alguma. O
termo $-\zeta_{\text{eq}}$ é o forçamento astronômico
(Capítulo~\ref{cap:potencial}, com o fator de Love \eqref{eq:love}); o fator
$(1-\beta)$ é a auto-atração e carga do Capítulo~\ref{cap:sal}.

\begin{table}[htbp]
\centering
\caption{Termos das Eq.~\eqref{eq:lte_cont}--\eqref{eq:lte_v}, o que cada um
representa e o valor adotado.}
\label{tab:termos}
\begin{tabular}{@{}llp{6.4cm}l@{}}
\toprule
termo & nome & papel físico & valor \\
\midrule
$f=2\Omega\sin\phi$ & Coriolis & sem ele não há anfidromo nem onda de Kelvin & --- \\
$H(\lambda,\phi)$ & batimetria & fixa $c=\sqrt{gH}$ e as fronteiras & ETOPO 2022 \\
$\zeta_{\text{eq}}$ & forçante & potencial de grau 2 (Lua $+$ Sol) & Eq.~\eqref{eq:heq} \\
$\beta$ & SAL & a maré deforma o geoide e a crosta & \num{0.09} \\
$C_d$ & atrito de fundo & camada limite turbulenta, $\sim\!70\%$ da dissipação & \num{0.0025} \\
$r$ & maré interna & conversão barotrópica $\to$ baroclínica & \SI{1e-5}{\per\second} \\
$\gamma_2$ & Terra elástica & $1+k_2-h_2$ & \num{0.69} \\
\bottomrule
\end{tabular}
\end{table}

\section{Condições de contorno}

\begin{itemize}
\item \textbf{Longitude periódica}: $\lambda=\ang{-180}$ e $\lambda=\ang{180}$
  são a mesma coluna de células, sem tratamento especial.
\item \textbf{Costas: fluxo normal nulo}, $\vect{u}\cdot\uvec{n}=0$. É a
  condição que \textbf{reflete} a onda — e portanto a que produz os sistemas
  anfidrômicos. Sem costas não há anfidromos; com costas eles são inevitáveis.
\item \textbf{Fronteira polar}: acima de \ang{86} de latitude tudo vira terra
  (\S\ref{sec:batimetria}).
\end{itemize}

\begin{observacao}[Por que não há condição de radiação]
Um domínio global fechado não tem fronteira aberta: a energia entra pelo
trabalho da forçante astronômica e sai pela dissipação. Isso torna o
Capítulo~\ref{cap:dissipacao} não opcional — sem termo dissipativo, o modelo
acumula energia indefinidamente, e é exatamente o que se mede na primeira linha
da Tabela~\ref{tab:atrito} ($r=0$ dá viés de $+\SI{33}{\centi\meter}$).
\end{observacao}

\section{Batimetria: a geometria como dado de entrada}
\label{sec:batimetria}

A maré de equilíbrio não precisa saber onde há água. Um modelo dinâmico precisa
de tudo, e essa é a segunda melhoria estrutural. A fonte é o \textbf{ETOPO 2022}
(\SI{1}{\arcminute}, NOAA NCEI), subamostrado no servidor via OPeNDAP para
\SI{5}{\arcminute} — divisor exato de \ang{0.25}, \ang{0.5} e \ang{1.0}, de
modo que qualquer resolução do modelo sai de um agrupamento inteiro de células.

Três decisões de tratamento do dado, todas com consequência física:

\begin{enumerate}
\item \textbf{A média de profundidade é feita apenas sobre as subcélulas
  oceânicas.} Misturar cotas positivas de terra na média rebaixaria
  artificialmente a profundidade das células costeiras — justamente onde a maré
  é mais sensível. Uma célula conta como água se a fração oceânica for
  $\ge\num{0.5}$.
\item \textbf{Profundidade mínima de \SI{20}{\meter}.} Em água muito rasa o
  atrito quadrático $C_d|\vect{u}|u/H$ diverge, e o esquema perde estabilidade.
  É o valor usual em modelos barotrópicos globais.
\item \textbf{Limite polar em \ang{86}.} A largura zonal da célula é
  $\Delta x=a\,\Delta\lambda\cos\phi$, que tende a zero no polo e levaria o
  passo estável junto (\S\ref{sec:cfl}). Em \ang{86} e \ang{1.0} de resolução a
  última linha de células molhadas, centrada em \ang{85.5}, ainda tem
  $\Delta x\approx\SI{8.7}{\kilo\meter}$, o que mantém
  $\Delta t\approx\SI{19}{\second}$. É restrição \emph{numérica}, e está
  declarada como tal — não é uma afirmação sobre a maré ártica.
\end{enumerate}

\section{Duas escolhas de método que definem a rodada}

\subsection{Forçamento monocromático}
\label{sec:monocromatico}

O modelo aceita dois regimes de forçamento, e o padrão de produção é o segundo:

\begin{description}[style=nextline,leftmargin=0.8cm]
\item[Astronômico] Lua e Sol a partir das efemérides, via \eqref{eq:heq}.
  Fisicamente mais rico, mas separar M\textsubscript{2} de S\textsubscript{2}
  na análise harmônica exige, pelo critério de Rayleigh
  (\S\ref{sec:rayleigh}), mais de \SI{14.8}{\day} de registro \emph{após} o
  spin-up.
\item[Monocromático] Uma constituinte de cada vez, na forma geodésica padrão de
  \S\ref{sec:especies}:
  \begin{equation}
  \zeta_{\text{eq}}=\gamma_2\,A_c\,S_{n}(\phi)\,
  \cos\!\left(\sigma_c t+n\lambda\right),
  \label{eq:forcante_mono}
  \end{equation}
  com $n$ a espécie e $S_n$ a forma latitudinal correspondente. A resposta é
  monocromática por construção, e bastam o spin-up e alguns períodos para
  extrair amplitude e fase com precisão. É como TPXO e FES operam.
\end{description}

O sinal $+n\lambda$ em \eqref{eq:forcante_mono} faz o padrão migrar para
\textbf{oeste} conforme o tempo avança, que é o sentido em que o ponto sublunar
se desloca visto de um referencial preso à Terra. Trocar esse sinal produz uma
solução espelhada, com anfidromos girando ao contrário — um erro que a
estatística de rotação do \S\ref{sec:rotacao} detectaria imediatamente.

\subsection{Análise harmônica embutida}
\label{sec:embutida}

Guardar o campo a cada passo custaria dezenas de gigabytes. Em vez disso, os
produtos $\zeta\cos\sigma t$ e $\zeta\sin\sigma t$ são acumulados
\textbf{durante} a integração e o sistema normal é resolvido no fim
(\S\ref{sec:minimosquadrados}). Guardam-se apenas os coeficientes; o campo em
qualquer instante é reconstruído deles em uma linha. É o mesmo princípio dos
modelos operacionais, e a diferença de custo é de três ordens de grandeza em
armazenamento.

\section{Energética}
\label{sec:energia}

Multiplicando \eqref{eq:lte_u}--\eqref{eq:lte_v} por $\rho H\vect{u}$,
\eqref{eq:lte_cont} por $\rho g(1-\beta)\zeta$, e somando, obtém-se a equação de
balanço
\begin{equation}
\pd{}{t}\underbrace{\left[\tfrac12\rho H|\vect{u}|^{2}
+\tfrac12\rho g(1-\beta)\zeta^{2}\right]}_{E}
+\nabla\cdot\underbrace{\left[\rho g(1-\beta)\,\zeta\,H\vect{u}\right]}_{\vect{F}\ \text{(fluxo)}}
=\underbrace{\rho g H\vect{u}\cdot\nabla\zeta_{\text{eq}}}_{W\ \text{(trabalho da maré)}}
-\underbrace{\rho\left(C_d|\vect{u}|^{3}+rH|\vect{u}|^{2}\right)}_{D\ \text{(dissipação)}} .
\label{eq:energia}
\end{equation}

O fluxo é o trabalho da pressão perturbada $p'=\rho g\zeta$ contra o transporte
$H\vect{u}$, e envolve $\zeta$ — não $\zeta'$. Escrevê-lo com $\zeta'$ seria
contabilizar o trabalho astronômico duas vezes, uma vez no divergente e outra no
termo $W$; a identidade deixaria de fechar. Com $\beta=0$ recupera-se a forma
usual, $E=\tfrac12\rho H|\vect{u}|^{2}+\tfrac12\rho g\zeta^{2}$ e
$\vect{F}=\rho g\zeta H\vect{u}$.

Note que o termo de Coriolis \textbf{não aparece}: ele é perpendicular à
velocidade e não realiza trabalho. Coriolis reorganiza a solução sem alimentá-la
— exatamente o que se espera de uma força inercial.

Em regime estatisticamente estacionário, $\langle W\rangle=\langle D\rangle$: a
potência astronômica injetada iguala a dissipada. A observação global dá
$\approx\SI{3.7}{\tera\watt}$ para todas as constituintes e
$\approx\SI{2.4}{\tera\watt}$ só para M\textsubscript{2}
\cite{egbert2000,munk1998}. A Eq.~\eqref{eq:energia} é o critério físico de
calibração do Capítulo~\ref{cap:dissipacao}: o coeficiente de dissipação não é
livre, ele precisa fechar esse orçamento.

% =====================================================================
\chapter{Auto-atração e carga (SAL)}
\label{cap:sal}

\section{O problema}

A maré é uma redistribuição de massa da ordem de $10^{16}\,\si{\kilo\gram}$
oscilando com período de \SI{12.42}{\hour} — a massa d'água acima do nível médio
integrada sobre o hemisfério em preamar. Essa massa faz três coisas que a
formulação do Capítulo~\ref{cap:lte} ignoraria se o termo $\beta$ não
existisse:

\begin{enumerate}
\item \textbf{atrai gravitacionalmente} o resto do oceano (auto-atração);
\item \textbf{deforma elasticamente} a crosta sobre a qual repousa (carga);
\item a deformação da crosta \textbf{redistribui massa sólida}, que por sua vez
  gera potencial adicional.
\end{enumerate}

Os três efeitos somam-se num termo que altera a equipotencial efetiva. Ignorá-lo
introduz erro de até \SI{20}{\percent} na amplitude global — não é uma
correção de segunda ordem.

\section{Formulação exata}

Expandindo a elevação da maré em harmônicos esféricos,
$\zeta=\sum_{n,m}\zeta_{nm}Y_{nm}$, o potencial de auto-atração e carga
correspondente, expresso como elevação equivalente, é
\begin{equation}
\boxed{\;
\zeta_{\text{SAL}}=\sum_{n}\frac{3\rho_w}{(2n+1)\rho_\oplus}
\left(1+k'_n-h'_n\right)\zeta_n,\;}
\label{eq:sal_exato}
\end{equation}
onde $\zeta_n$ é a componente de grau $n$ do campo de maré, $\rho_w/\rho_\oplus
=\num{1025}/\num{5515}=\num{0.186}$, e $k'_n,h'_n$ são os \emph{números de Love
de carga} — negativos, ao contrário dos números de Love de maré do
\S\ref{sec:love}, porque uma carga superficial \emph{afunda} a crosta:
$k'_2\approx\num{-0.31}$, $h'_2\approx\num{-1.00}$.

O termo de SAL entra na Eq.~\eqref{eq:desvio} como
$\zeta'=\zeta-\zeta_{\text{eq}}-\zeta_{\text{SAL}}$.

\section{A aproximação escalar e o seu erro}

A Eq.~\eqref{eq:sal_exato} é uma \textbf{convolução global}: cada ponto do
oceano depende de todos os outros, a cada passo de tempo. Calculá-la
exigiria uma transformada de harmônicos esféricos por passo, o que multiplica
o custo do modelo por um fator grande.

A aproximação padrão, devida a Accad \& Pekeris \cite{accad1978} e avaliada em
detalhe por Ray \cite{ray1998}, é escalar:
\begin{equation}
\zeta_{\text{SAL}}\approx\beta\,\zeta,
\qquad\beta\in[\num{0.085},\,\num{0.12}],
\end{equation}
o que produz exatamente a Eq.~\eqref{eq:desvio}. O valor adotado aqui é
$\beta=\num{0.09}$.

\begin{caixa}[De onde vem o número \num{0.09}]
Avalie o coeficiente de \eqref{eq:sal_exato} grau a grau:
\begin{equation}
\beta_n=\frac{3\rho_w}{(2n+1)\rho_\oplus}\left(1+k'_n-h'_n\right)
\;\Longrightarrow\;
\beta_2\approx\num{0.19},\quad
\beta_{10}\approx\num{0.06},\quad
\beta_n\xrightarrow[n\to\infty]{}0 .
\end{equation}
O escalar $\beta$ não é nenhum $\beta_n$: é uma \textbf{média ponderada pelo
espectro espacial da própria maré}. Como a maré real tem energia distribuída
por muitos graus — e não concentrada em $n=2$ como a forçante —, o valor
efetivo cai para $\sim\!\num{0.09}$. Isso também explica por que a aproximação
escalar erra mais em plataforma continental, onde o espectro da maré é
deslocado para graus altos: ali o erro chega a \SI{20}{\percent}.
\end{caixa}

\begin{observacao}[Como o erro se manifesta]
Por ser proporcional a $\zeta$, o termo escalar age como uma redução efetiva de
$g$: $g\to g(1-\beta)$ na Eq.~\eqref{eq:lte_u}. Isso reduz a velocidade de fase
em $\sqrt{1-\beta}\approx\num{0.954}$, ou seja, cerca de \SI{4.6}{\percent}, e
desloca ligeiramente todos os modos ressonantes da bacia. É um efeito de
\emph{fase} tanto quanto de amplitude — motivo pelo qual modelos operacionais
que buscam \SI{1}{\centi\meter} de acurácia abandonaram a aproximação escalar.
\end{observacao}

% =====================================================================
\chapter{Dissipação e orçamento energético}
\label{cap:dissipacao}

Num domínio global fechado, a dissipação não é um detalhe de acabamento: é o
que fixa a amplitude. Sem ela, \eqref{eq:energia} garante crescimento
indefinido enquanto houver trabalho da forçante. Este capítulo formula os dois
mecanismos implementados e a metodologia de calibração do único parâmetro
verdadeiramente livre do modelo.

\section{Atrito de fundo quadrático}

Na camada limite turbulenta junto ao fundo, a tensão de cisalhamento escala com
o quadrado da velocidade — resultado clássico de G.~I. Taylor para o Mar da
Irlanda \cite{taylor1919}:
\begin{equation}
\vect{\tau}_b=\rho\,C_d\,|\vect{u}|\,\vect{u},
\qquad C_d\approx\num{2.5e-3}.
\end{equation}
Distribuída sobre a coluna d'água de espessura $H$, a desaceleração
correspondente é $C_d|\vect{u}|\vect{u}/H$, que é o termo que aparece em
\eqref{eq:lte_u}--\eqref{eq:lte_v}. A dependência $1/H$ concentra a dissipação
onde a água é rasa: mares marginais e plataformas continentais respondem por
cerca de \SI{70}{\percent} do total, embora ocupem uma fração pequena da área
oceânica.

A dissipação instantânea por unidade de área é $\rho C_d|\vect{u}|^{3}$: cúbica
na velocidade. Isso torna o termo fortemente seletivo — dobrar a corrente
octuplica a perda —, e é a razão pela qual ele estabiliza o modelo sem precisar
de ajuste fino.

\section{Conversão para maré interna}

O segundo mecanismo é invisível a um modelo barotrópico. Quando a corrente de
maré atravessa topografia rugosa num oceano estratificado, parte da energia é
convertida em \textbf{ondas internas} que se propagam ao longo da picnoclina e
acabam quebrando longe dali. Egbert \& Ray \cite{egbert2000}, usando altimetria
do TOPEX/Poseidon, atribuíram $\approx\SI{1}{\tera\watt}$ — 25 a
\SI{30}{\percent} do total — a esse sumidouro de oceano profundo, resultado que
reescreveu o orçamento energético das marés.

Um modelo de densidade constante não pode representar essas ondas
(\S\ref{sec:aguasrasas}). Pode apenas \emph{parametrizar} a perda que elas
representam. A forma usada aqui é um arrasto linear
\begin{equation}
\vect{\mathcal{F}}_{\text{int}}=-r\,\vect{u},
\qquad r=\SI{1e-5}{\per\second},
\label{eq:arrastolinear}
\end{equation}
uniforme. Parametrizações mais sofisticadas tornam $r$ função da rugosidade e
da estratificação locais — a de Jayne \& St.~Laurent \cite{jayne2001} usa
$r\propto\kappa h_t^{2}N_b/H$ com $h_t$ a altura característica da topografia e
$N_b$ a frequência de Brunt--Väisälä junto ao fundo. Aqui, deliberadamente,
mantém-se o escalar: com um único parâmetro, a calibração é auditável.

\begin{caixa}[Por que \emph{linear} e não quadrático]
A escolha da forma funcional não é arbitrária. O atrito de fundo é quadrático
porque a camada limite é turbulenta. A conversão para maré interna, ao
contrário, é um processo de \emph{radiação de ondas}: a energia irradiada é
proporcional ao quadrado da amplitude da corrente, ou seja, a força equivalente
é linear na velocidade. As duas formas funcionais codificam mecanismos físicos
distintos, e usar a mesma para ambos apagaria a distinção.
\end{caixa}

\section{Calibração: metodologia}
\label{sec:calibracao}

O coeficiente $r$ \textbf{não é derivável do modelo}. É um parâmetro fechado
contra observação, e a metodologia importa mais que o valor.

O procedimento: varrer $r$, e para cada valor rodar o modelo completo a
\ang{1.0} com forçante monocromática de M\textsubscript{2}, spin-up de
\SI{10}{\day} e análise de \SI{2.5}{\day}, comparando a amplitude resultante
com as constantes harmônicas de \textbf{oito marégrafos insulares} da NOAA
(a justificativa da escolha das estações está em \S\ref{sec:porqueilhas}).

\begin{table}[htbp]
\centering
\caption{Varredura de calibração do arrasto de maré interna. O mínimo é
limpo e, decisivamente, o
\emph{viés troca de sinal} exatamente nele.}
\label{tab:atrito}
\begin{tabular}{@{}rrrrr@{}}
\toprule
$r$ [\si{\per\second}] & RMS ilhas [cm] & viés [cm] & máx.\ global [m] & tempo [s] \\
\midrule
\num{0}      & \num{36.13} & $+\num{32.89}$ & \num{7.03} & 163 \\
\num{3e-6}   & \num{18.58} & $+\num{16.00}$ & \num{4.69} & 172 \\
\textbf{\num{1e-5}} & \textbf{\num{5.24}} & $+\num{0.13}$ & \num{2.91} & 169 \\
\num{2e-5}   & \num{9.17}  & $-\num{8.04}$  & \num{2.10} & 167 \\
\num{3e-5}   & \num{13.07} & $-\num{11.87}$ & \num{1.59} & 167 \\
\num{5e-5}   & \num{17.36} & $-\num{15.60}$ & \num{0.95} & 167 \\
\bottomrule
\end{tabular}
\end{table}

\begin{caixa}[O que torna esta calibração defensável]
Um mínimo de RMS, isoladamente, pode ser artefato de compensação entre erros de
sinais opostos. O que valida a Tabela~\ref{tab:atrito} é a
\textbf{coluna do viés}: ela é monótona e cruza zero exatamente onde o RMS é
mínimo. Isso significa que em $r=\SI{1e-5}{\per\second}$ o modelo não está
acertando na média por cancelamento — está centrado. Além disso, a amplitude
máxima global decresce monotonicamente e passa por valores fisicamente
plausíveis ($\sim\SI{3}{\meter}$) justamente ali, enquanto $r=0$ produz
\SI{7}{\meter}, incompatível com qualquer observação.
\end{caixa}

\begin{alerta}[Um erro de calibração e como foi detectado]
A primeira versão usava $r=\SI{2e-4}{\per\second}$, vinte vezes o valor final.
O diagnóstico veio de uma checagem de escala de tempo, não de um ajuste: o
tempo de decaimento associado é $\tau=1/r=\SI{1.4}{\hour}$, contra um período
de M\textsubscript{2} de \SI{12.42}{\hour}. Um oscilador cujo transiente decai
nove vezes mais rápido do que o período do forçamento está
\emph{superamortecido} — a resposta não pode ressoar, e nenhuma bacia
amplifica. Comparar a constante de tempo do amortecimento com o período do
forçamento é um teste de sanidade que se faz antes de qualquer varredura.
\end{alerta}

\section{Orçamento global e o torque de maré}
\label{sec:torque}

A dissipação de maré não é apenas um termo de fechamento numérico: é um dos
números mais bem determinados da geofísica, e conecta este modelo à mecânica
celeste do Capítulo~\ref{cap:gravitacao}.

\begin{center}
\begin{tabular}{@{}lrl@{}}
\toprule
sumidouro & potência & referência \\
\midrule
total, todas as constituintes & \SI{3.7}{\tera\watt} & \cite{munk1998} \\
M\textsubscript{2} apenas & \SI{2.4}{\tera\watt} & \cite{egbert2000} \\
mares rasos e plataformas & $\sim\SI{2.6}{\tera\watt}$ & atrito de fundo \\
oceano profundo (maré interna) & $\sim\SI{1.0}{\tera\watt}$ & \cite{egbert2000} \\
\bottomrule
\end{tabular}
\end{center}

O mecanismo de fundo é o seguinte. Por causa da dissipação, o bojo de maré não
está alinhado com o eixo Terra--Lua: está adiantado de um pequeno ângulo
(\S\ref{sec:deslocamento}). Um bojo desalinhado exerce \textbf{torque} sobre a
Lua, e a reação frena a rotação da Terra. As consequências são medidas com
precisão:

\begin{itemize}
\item a Lua afasta-se a \SI{3.8}{\centi\meter\per\ano} (telemetria laser
  lunar);
\item o dia alonga-se $\sim\SI{2.3}{\milli\second}$ por século;
\item o momento angular perdido pela rotação terrestre reaparece
  \emph{integralmente} na órbita lunar. A energia, essa não se conserva: a
  energia rotacional perdida excede o ganho de energia orbital, e a diferença é
  dissipada como calor — os \SI{3.7}{\tera\watt}.
\end{itemize}

\begin{caixa}[O fecho conceitual]
Este é o ponto em que todo o documento se articula. O termo de arrasto do
Capítulo~\ref{cap:referenciais} gerou a força; a expansão multipolar do
Capítulo~\ref{cap:gravitacao} isolou o quadrupolo; a hidrodinâmica do
Capítulo~\ref{cap:lte} produziu uma resposta \emph{atrasada}; e o atraso, por
sua vez, realimenta a mecânica celeste que iniciou o problema — alterando a
órbita da Lua e a rotação da Terra. A maré não é um efeito passivo da órbita:
é um termo do próprio problema de dois corpos, e é ele que torna o sistema
Terra--Lua dissipativo.
\end{caixa}

% =====================================================================
\chapter{Coriolis, ondas de Kelvin e pontos anfidrômicos}
\label{cap:anfidromos}

Este capítulo trata do fenômeno que a teoria de equilíbrio é estruturalmente
incapaz de produzir, e que constitui a validação qualitativa mais forte do
modelo dinâmico.

\section{Do problema de Taylor ao anfidromo}

Considere um canal semi-infinito, girante, fechado numa extremidade — a
idealização de uma bacia oceânica. Uma onda de Kelvin \eqref{eq:kelvin} incide
propagando-se ao longo de uma margem; ao encontrar a parede, ela não pode
simplesmente refletir como numa cuba sem rotação, porque a onda de Kelvin só
existe com a costa à direita (hemisfério norte). A solução, obtida por Taylor
em 1922 \cite{taylor1922}, é que a energia refletida sai como uma segunda onda
de Kelvin \emph{ao longo da margem oposta}, mais um conjunto de modos de
Poincaré evanescentes que ajustam o campo perto da parede.

A superposição de duas ondas de Kelvin, propagando-se em sentidos opostos
junto a margens opostas e decaindo transversalmente, produz um campo cuja
amplitude se anula em pontos isolados do interior do canal. Em torno desses
pontos, a fase gira. São os \textbf{pontos anfidrômicos}.

\begin{figure}[htbp]
\centering
\begin{tikzpicture}[scale=0.95,>=Latex]
  \fill[azulclaro] (0,0) rectangle (8,5);
  \draw[very thick,azulforte] (0,0) -- (8,0);
  \draw[very thick,azulforte] (0,5) -- (8,5);
  \draw[very thick,azulforte] (8,0) -- (8,5);
  \node[below=2pt,font=\small] at (4,0) {costa};
  \node[above=2pt,font=\small] at (4,5) {costa};
  \node[right=3pt,font=\small] at (8,2.5) {parede};
  % linhas cotidais radiais, com folga vertical em 270 graus
  \foreach \a in {15,45,...,345}{
    \draw[azulforte,thin] (5.4,2.5) -- ++(\a:1.5);
  }
  \fill[laranja] (5.4,2.5) circle (3pt);
  \draw[gray!70,thin] (5.4,1.02) -- (5.4,0.62);
  \node[below=-1pt,font=\small] at (5.4,0.68) {anfidromo};
  % sentido de rotacao
  \draw[->,laranja,very thick] (5.4,2.5)+(-45:1.95) arc (-45:45:1.95);
  % ondas de Kelvin
  \draw[->,verde,very thick] (0.4,0.45) -- (2.8,0.45);
  \node[above=1pt,font=\scriptsize,verde] at (1.6,0.5) {Kelvin incidente};
  \draw[->,verde,very thick] (2.8,4.55) -- (0.4,4.55);
  \node[below=1pt,font=\scriptsize,verde] at (1.6,4.5) {Kelvin refletida};
\end{tikzpicture}
\caption{Problema de Taylor: a reflexão de uma onda de Kelvin no extremo
fechado de um canal girante produz um ponto anfidrômico, em torno do qual as
linhas cotidais (de igual fase) convergem em leque. A amplitude é nula no
ponto e cresce radialmente; a fase percorre \ang{360} ao circundá-lo. A seta
em laranja marca o sentido de propagação da fase, anti-horário no hemisfério
norte.}
\label{fig:taylor}
\end{figure}

\section{O anfidromo como objeto topológico}
\label{sec:topologia}

Este é o ponto de vista que torna a detecção rigorosa. Escreva a resposta
monocromática como a parte real de um campo complexo:
\begin{equation}
\zeta(\vect{x},t)=\mathrm{Re}\left\{Z(\vect{x})\,e^{\ii\sigma t}\right\}
=A\cos(\sigma t-G),
\qquad Z=A(\vect{x})\,e^{-\ii G(\vect{x})},
\label{eq:complexo}
\end{equation}
com $A\ge 0$ a amplitude e $G$ a fase de Greenwich — o \emph{atraso} com que a
preamar chega, na mesma convenção de \eqref{eq:previsao} e de
\S\ref{sec:minimosquadrados}, onde $G=\atantwo(b,a)$. A convenção importa: com
$\zeta=A\cos(\sigma t+G)$ o sinal do número de voltas \eqref{eq:winding}
inverteria, e com ele a leitura de sentido de rotação da
Tabela~\ref{tab:rotacao}. Um ponto anfidrômico é
simplesmente um \textbf{zero de $Z$}.

Como $Z:\mathbb{R}^{2}\to\mathbb{C}$ é uma aplicação de duas variáveis reais em
duas, seus zeros são genericamente \emph{isolados} — anular simultaneamente
$\mathrm{Re}\,Z$ e $\mathrm{Im}\,Z$ são duas condições em duas incógnitas. E a
cada zero isolado associa-se um \textbf{índice topológico}, o número de voltas
da fase:
\begin{equation}
\boxed{\;
w=\frac{1}{2\pi}\oint_{\mathcal{C}}\nabla G\cdot\dd\vect{l}\;\in\;\mathbb{Z},\;}
\label{eq:winding}
\end{equation}
para qualquer curva fechada $\mathcal{C}$ que circunde o ponto e não contenha
outros zeros. Genericamente $w=\pm 1$, e o \textbf{sinal dá o sentido de
rotação} da onda em torno do ponto.

\begin{caixa}[Por que anfidromos são robustos]
O índice \eqref{eq:winding} é um invariante topológico: perturbe a batimetria,
o atrito ou a forçante continuamente, e $w$ não pode mudar sem que o zero
colida com outro de índice oposto e ambos se aniquilem — ou que ele saia pela
fronteira do domínio, que é como surgem os anfidromos \emph{degenerados}, com
centro em terra (o do Mar do Norte junto à costa da Noruega é o exemplo
padrão). Isso explica um fato
observacional notável — os anfidromos oceânicos estão nos mesmos lugares em
todos os modelos e em todas as épocas, apesar de a solução detalhada ser
sensível a tudo. Eles não são ``pontos onde a maré é pequena''; são
\textbf{defeitos topológicos} do campo de fase, e é por isso que a detecção
correta precisa ser topológica.
\end{caixa}

\section{Algoritmo de detecção}

A tradução direta de \eqref{eq:winding} para a grade discreta: percorra o anel
dos oito vizinhos de cada célula em ordem angular e some as diferenças de fase,
cada uma reduzida ao intervalo $(-\ang{180},\ang{180}]$:
\begin{equation}
W_{ji}=\sum_{k=1}^{8}
\Big[\big(G_{k+1}-G_{k}+\ang{180}\big)\bmod\ang{360}-\ang{180}\Big].
\end{equation}
O total vale $\pm\ang{360}$ num anfidromo verdadeiro e \ang{0} em qualquer
outro lugar. Detecções vizinhas são fundidas por proximidade, com o
representante do grupo sendo a célula de menor amplitude.

A tradução direta para código está no Apêndice~\ref{ap:python}.

\begin{alerta}[O critério ingênuo e por que ele falha]
A primeira implementação detectava anfidromos por ``amplitude abaixo de um
limiar'' e encontrava \textbf{86} pontos numa rodada a \ang{1.0} — muitos deles
em baías fechadas e em ruído costeiro, onde a maré é pequena por razões que
nada têm de anfidrômicas. Com o critério \eqref{eq:winding} mais o
agrupamento por proximidade, o número cai para \textbf{24}, compatível com as
cartas cotidais publicadas. A diferença entre \num{86} e \num{24} não é ajuste
de limiar: é a diferença entre uma definição operacional imprecisa e a
definição correta.
\end{alerta}

\section{O sentido de rotação emerge de $f$}
\label{sec:rotacao}

A onda de Kelvin costeira circula com a costa à direita no hemisfério norte, o
que faz os sistemas anfidrômicos girarem majoritariamente no sentido
\textbf{anti-horário} ao norte e \textbf{horário} ao sul. Nada disso está
codificado; sai do termo $f=2\Omega\sin\phi$.

\begin{table}[htbp]
\centering
\caption{Sentido de rotação medido na rodada de produção a \ang{1.0}, que
encontra 24 pontos anfidrômicos.}
\label{tab:rotacao}
\begin{tabular}{@{}lrr@{}}
\toprule
hemisfério & anfidromos & no sentido previsto \\
\midrule
norte (anti-horário) & 8 & \SI{75}{\percent} \\
sul (horário) & 13 & \SI{85}{\percent} \\
\midrule
equatoriais ($|\phi|<\ang{5}$) & 3 & não classificados \\
\bottomrule
\end{tabular}
\end{table}

\noindent
Os três pontos dentro de \ang{5} do equador ficam de fora da estatística
porque ali $f\to 0$: não há sentido preferencial a prever, e incluí-los
contaminaria o teste com casos em que a hipótese não se aplica.

\begin{figure}[htbp]
\centering
\includegraphics[width=\textwidth]{carta_cotidal_M2.png}
\caption{Carta cotidal de M\textsubscript{2} produzida pelo modelo a \ang{1.0}.
Cor: amplitude (escala sequencial de um só tom, porque amplitude é magnitude e
não polaridade). Linhas: fase constante — cada uma é uma frente de onda. Os
pontos onde todas as linhas convergem e a amplitude vai a zero são os
anfidromos. A carta reproduz os sistemas reais conhecidos: Atlântico Sul em
$\sim(\ang{-25},\ang{-12})$, Atlântico Norte em $\sim(\ang{50},\ang{-40})$,
Índico em $\sim(\ang{0},\ang{62})$.}
\label{fig:cotidal}
\end{figure}

\section{No anfidromo a maré não sobe, mas a água corre}

Uma verificação que separa compreensão de repetição. No ponto anfidrômico a
elevação é nula por construção; a \emph{corrente}, não:

\begin{center}
\begin{tabular}{@{}lr@{}}
\toprule
grandeza no anfidromo & fração da mediana global \\
\midrule
elevação & \SI{6.4}{\percent} \\
corrente & \SI{102}{\percent} \\
\bottomrule
\end{tabular}
\end{center}

Como $Z=0$ no anfidromo, ali $\zeta\equiv 0$ em \emph{todo} instante, e com ela
o fluxo de energia $\rho g\zeta H\vect{u}$ de \eqref{eq:energia}: a energia
\textbf{contorna} o ponto, não o atravessa — nas cartas de fluxo de maré os
vetores circulam em torno do anfidromo com magnitude indo a zero nele. O que
não se anula é a \textbf{corrente}: o transporte de volume $H\vect{u}$ passa
normalmente. Um anfidromo é um ponto de maré nula, não de água parada.

\begin{alerta}[Uma formulação que se erra com facilidade]
É tentador resumir o parágrafo acima dizendo que ``elevação e corrente estão em
quadratura''. Isso vale para a \emph{fase}, não para a \emph{amplitude}.
Espacialmente as duas amplitudes são bem correlacionadas
($\rho\approx\num{0.79}$), porque plataforma rasa tem maré grande \emph{e}
corrente grande. A afirmação que a física sustenta é pontual — sobre o
anfidromo —, e não uma anticorrelação global entre os dois campos.
\end{alerta}

% =====================================================================
\chapter{Análise harmônica}
\label{cap:harmonica}

Há dois caminhos para prever maré, e eles são filosoficamente distintos:
\textbf{simular} a hidrodinâmica, ou \textbf{decompor} o sinal em constituintes
de frequência astronomicamente determinada. O segundo caminho é o que produz as
tábuas de maré do mundo inteiro, e é um instrumento tão essencial quanto as
próprias Equações de Maré de Laplace — porque é ele que torna a simulação
\emph{comparável} com a observação.

\section{A ideia de Kelvin e Darwin}

A sacada é a seguinte: as \emph{frequências} da maré saem da mecânica celeste
com precisão arbitrária, ao passo que amplitude e fase dependem da
hidrodinâmica local, impossível de calcular no século XIX. Logo, fixem-se as
frequências pela astronomia e \emph{meça-se} o resto:
\begin{equation}
\zeta(t)=Z_0+\sum_{i}f_i\,H_i\,\cos\!\left(V_{0i}+u_i-g_i\right),
\label{eq:previsao}
\end{equation}
onde $H_i$ e $g_i$ (a \emph{fase de Greenwich}) são as constantes harmônicas
medidas naquele porto, $V_{0i}$ é o argumento astronômico e $f_i,u_i$ são as
correções nodais. A Eq.~\eqref{eq:previsao} é literalmente o algoritmo de uma
tábua de maré.

\section{Os seis argumentos fundamentais e os números de Doodson}

Toda constituinte é uma combinação com coeficientes \emph{inteiros} de seis
frequências fundamentais:

\begin{table}[htbp]
\centering
\caption{Os seis argumentos fundamentais de Doodson, em graus por hora solar
média.}
\label{tab:doodson}
\begin{tabular}{@{}cllr@{}}
\toprule
símbolo & significado & período & \si{\degree\per\hour} \\
\midrule
$\tau$ & tempo lunar médio & \SI{24.84}{\hour} & \num{14.4920521} \\
$s$ & longitude média da Lua & \SI{27.32}{\day} & \num{0.5490165} \\
$h$ & longitude média do Sol & \SI{365.24}{\day} & \num{0.0410686} \\
$p$ & perigeu lunar & \SI{8.85}{\anos} & \num{0.0046418} \\
$N'$ & nó lunar (negativo) & \SI{18.61}{\anos} & \num{0.0022064} \\
$p_s$ & perigeu solar & \SI{20940}{\anos} & \num{0.0000020} \\
\bottomrule
\end{tabular}
\end{table}

A velocidade angular de uma constituinte é o produto escalar do seu vetor de
Doodson $(n_\tau,n_s,n_h,n_p,n_{N},n_{p_s})$ com a coluna da direita:
\begin{equation}
\sigma_c=\sum_{k}n_k\,\omega_k .
\label{eq:velocidade}
\end{equation}

Assim, M\textsubscript{2} é $(2,0,0,0,0,0)$, logo
$2\times\num{14.4920521}=\SI{28.9841042}{\degree\per\hour}$, ou seja
\SI{12.4206}{\hour} de período — exatamente o valor que a simulação obtém por
FFT sem jamais tê-lo codificado. E o primeiro número de Doodson $n_\tau$
\emph{é} a espécie de \S\ref{sec:especies}: 0 para longo período, 1 para
diurna, 2 para semidiurna.

A relação de Doodson que fecha o sistema é
\begin{equation}
\tau=T+h-s,\qquad T=15\times(\text{hora UT})+\ang{180},
\end{equation}
com $T$ o ângulo horário médio do \emph{Sol}. É essa subtração que dá o ``2''
de M\textsubscript{2} referido ao dia \emph{lunar}, e não ao solar.

\begin{table}[htbp]
\centering
\caption{As oito constituintes principais, contra as quais a acurácia dos
modelos globais é medida \cite{stammer2014}. Amplitudes de equilíbrio do
desenvolvimento de Cartwright--Tayler--Edden.}
\label{tab:constituintes}
\begin{tabular}{@{}llrrl@{}}
\toprule
nome & Doodson & \si{\degree\per\hour} & $A$ [m] & descrição \\
\midrule
M\textsubscript{2} & $(2,0,0,0,0,0)$ & \num{28.9841} & \num{0.2423} & principal lunar semidiurna \\
S\textsubscript{2} & $(2,2,-2,0,0,0)$ & \num{30.0000} & \num{0.1128} & principal solar semidiurna \\
N\textsubscript{2} & $(2,-1,0,1,0,0)$ & \num{28.4397} & \num{0.0464} & elíptica lunar maior \\
K\textsubscript{2} & $(2,2,0,0,0,0)$ & \num{30.0821} & \num{0.0307} & declinacional luni-solar \\
K\textsubscript{1} & $(1,1,0,0,0,0)$ & \num{15.0411} & \num{0.1416} & declinacional luni-solar \\
O\textsubscript{1} & $(1,-1,0,0,0,0)$ & \num{13.9430} & \num{0.1005} & principal lunar diurna \\
P\textsubscript{1} & $(1,1,-2,0,0,0)$ & \num{14.9589} & \num{0.0468} & principal solar diurna \\
Q\textsubscript{1} & $(1,-2,0,1,0,0)$ & \num{13.3987} & \num{0.0193} & elíptica lunar diurna \\
\bottomrule
\end{tabular}
\end{table}

\section{Correções de fase por espécie}
\label{sec:correcaofase}

O argumento astronômico completo de uma constituinte é
\begin{equation}
V_0=\sum_k n_k\,\theta_k+\chi_c ,
\end{equation}
onde $\chi_c$ é uma \textbf{correção de fase constante} que depende apenas da
espécie. Ela decorre da convenção com que o argumento de Doodson é tabulado: a
origem de fase adotada para cada espécie no desenvolvimento harmônico não
coincide com a de \eqref{eq:especies}, e o descompasso é de \ang{0}, \ang{90} e
\ang{180} para as espécies 2, 1 e 0, respectivamente. O sinal, dentro da espécie
1, acompanha o sinal do coeficiente da constituinte no desenvolvimento.

\begin{center}
\small
\begin{tabular}{@{}llp{8.2cm}@{}}
\toprule
espécie & $\chi_c$ & origem \\
\midrule
semidiurnas & \ang{0} & a origem de fase do argumento coincide com a de \eqref{eq:especies} \\
diurnas & $\pm\ang{90}$ & origem deslocada de \ang{90} para a espécie 1; K\textsubscript{1} leva sinal oposto ao de O\textsubscript{1}, P\textsubscript{1} e Q\textsubscript{1} \\
longo período & \ang{180} & origem deslocada de \ang{180}, coerente com $\tfrac12(3\sin^{2}\phi-1)<0$ em baixa latitude \\
\bottomrule
\end{tabular}
\end{center}

\begin{alerta}[O erro de \ang{180} e o experimento que o revelou]
Na primeira implementação todas as correções diurnas estavam com o
\textbf{sinal invertido}. O erro é invisível numa inspeção de código — a tabela
parecia razoável — e foi encontrado por um experimento independente: ajustar as
\emph{próprias previsões oficiais da NOAA} com os nossos argumentos
astronômicos e verificar se recuperávamos as constantes que a NOAA publica.

O resultado foi diagnóstico: as semidiurnas saíram exatas
($\Delta g_{M2}=\ang{0.0}$) enquanto \emph{todas} as diurnas erravam por
exatamente $\pm\ang{180}$. Um erro que atinge uma espécie inteira e é um
múltiplo exato de $\pi$ não é ruído de ajuste — é sinal trocado. Corrigido o
sinal, a correlação da reconstrução subiu de \num{0.40} para acima de
\num{0.99}.

A lição metodológica: \textbf{teste contra uma fonte externa que use a mesma
convenção}, e prefira métricas que revelem estrutura (aqui, o erro por espécie)
a métricas agregadas (um RMS global teria apenas ficado ``ruim'').
\end{alerta}

\section{Fatores nodais}

A órbita lunar precessa com período de \SI{18.61}{\anos} (\S\ref{sec:multipolo}),
o que modula amplitude e fase de cada constituinte. Em vez de introduzir
constituintes separadas para cada linha espectral vizinha, a prática padrão —
devida a Schureman — é escrever
\begin{equation}
H_i\to f_i(N)\,H_i,\qquad g_i\to g_i-u_i(N),
\end{equation}
com $f$ e $u$ funções lentas da longitude do nó $N$. Alguns exemplos:
\begin{align}
f_{M2}&=1-\num{0.037}\cos N, & u_{M2}&=-\ang{2.14}\sin N,\\
f_{K1}&=\num{1.006}+\num{0.115}\cos N-\num{0.009}\cos 2N,
& u_{K1}&=-\ang{8.86}\sin N+\ang{0.68}\sin 2N,\\
f_{O1}&=\num{1.009}+\num{0.187}\cos N-\num{0.015}\cos 2N,
& u_{O1}&=\ang{10.80}\sin N-\ang{1.34}\sin 2N .
\end{align}

A modulação é de $\pm\SI{3.7}{\percent}$ para M\textsubscript{2} mas
$\pm\SI{18}{\percent}$ para O\textsubscript{1}: as diurnas são muito mais
sensíveis ao ciclo nodal, porque dependem de $\sin 2\delta$
(\S\ref{sec:especies}) e a declinação máxima da Lua oscila entre \ang{18.3} e
\ang{28.6}. Sem essa correção, uma análise harmônica de um ano produz
constantes que não valem para o ano seguinte.

As constituintes de águas rasas herdam o fator do harmônico gerador elevado à
potência apropriada: $f_{M4}=f_{M2}^{2}$, $f_{M6}=f_{M2}^{3}$,
$u_{M4}=2u_{M2}$.

\section{Critério de Rayleigh e condicionamento}
\label{sec:rayleigh}

Quanto tempo de registro é preciso para separar duas constituintes? O
\textbf{critério de Rayleigh} exige ao menos um ciclo completo do batimento:
\begin{equation}
T_{\text{reg}}\;\ge\;\frac{2\pi}{|\sigma_1-\sigma_2|} .
\label{eq:rayleigh}
\end{equation}

\begin{center}
\begin{tabular}{@{}llr@{}}
\toprule
par & $\Delta\sigma$ [\si{\degree\per\hour}] & registro mínimo \\
\midrule
M\textsubscript{2}, S\textsubscript{2} & \num{1.0159} & \SI{14.77}{\day} \\
K\textsubscript{1}, O\textsubscript{1} & \num{1.0981} & \SI{13.66}{\day} \\
M\textsubscript{2}, N\textsubscript{2} & \num{0.5444} & \SI{27.55}{\day} \\
K\textsubscript{1}, P\textsubscript{1} & \num{0.0821} & \textbf{\SI{182.6}{\day}} \\
\bottomrule
\end{tabular}
\end{center}

O par K\textsubscript{1}/P\textsubscript{1} exige meio ano — motivo pelo qual
registros curtos inferem P\textsubscript{1} a partir de K\textsubscript{1} por
razão de amplitude teórica, em vez de ajustá-la.

\begin{caixa}[A generalização quantitativa: número de condição]
O critério \eqref{eq:rayleigh} é binário e, estritamente, não é uma barreira
absoluta: com dados sem ruído, mínimos quadrados separam frequências
arbitrariamente próximas. O que \eqref{eq:rayleigh} realmente marca é a
transição de regime do \textbf{condicionamento} da matriz normal
$\vect{N}=\vect{X}^{\!\top}\vect{X}$ de \S\ref{sec:minimosquadrados}.

Medindo $\mathrm{cond}(\vect{N})$ em função do comprimento do registro para o
par K\textsubscript{1}/P\textsubscript{1}, encontra-se decrescimento monótono
até atingir \num{1.00} — a duas casas decimais — no período de Rayleigh de
\SI{183}{\day}, onde as colunas de projeto se tornam ortogonais a menos de um
resíduo $\ordem{\Delta\sigma/(\sigma_1+\sigma_2)}$, da ordem de
$\num{e-3}$. Abaixo disso o
condicionamento explode, e o efeito prático é que a variância dos coeficientes
estimados, $\sigma^{2}\vect{N}^{-1}$, cresce sem limite. A verificação honesta
do critério mede essa monotonicidade e o erro sobre realizações ruidosas, em vez
de tratar \eqref{eq:rayleigh} como dogma.
\end{caixa}

\section{Mínimos quadrados e acumulação em tempo de integração}
\label{sec:minimosquadrados}

Com as frequências fixadas, \eqref{eq:previsao} é \textbf{linear} nos
parâmetros. Escrevendo cada constituinte como
$a_i\cos\sigma_i t+b_i\sin\sigma_i t$, a solução de mínimos quadrados é
\begin{equation}
\vect{N}\vect{c}=\vect{y},\qquad
N_{kl}=\sum_{t}\Psi_k(t)\Psi_l(t),\qquad
y_k=\sum_{t}\Psi_k(t)\,\zeta(t),
\label{eq:normal}
\end{equation}
com $\Psi=(\cos\sigma_1 t,\sin\sigma_1 t,\dots)$. Amplitude e fase saem de
$A=\sqrt{a^{2}+b^{2}}$ e $G=\atantwo(b,a)$.

O ponto de implementação decisivo é que \eqref{eq:normal} envolve apenas
\textbf{somas}: nada exige que a série temporal seja armazenada. O modelo
acumula $\vect{N}$ e $\vect{y}$ durante a integração e resolve o sistema no
fim, para elevação \emph{e} para as duas componentes de corrente. Como
$\vect{N}$ é a mesma para todas as células da grade, uma única fatoração
resolve os \num{65000} pontos de uma vez.



\section{Constituintes de águas rasas: geradas, não forçadas}

M\textsubscript{4}, M\textsubscript{6}, MS\textsubscript{4} e companhia têm
amplitude de equilíbrio \textbf{exatamente zero}: não existem no potencial
astronômico. Elas nascem das não-linearidades — atrito quadrático, o termo
$(H+\zeta)$ da continuidade, advecção — em mares rasos, onde
$\mathrm{Ro}\sim 1$ (\S\ref{sec:escala}).

Essa é uma distinção conceitual de primeira importância, e um teste
discriminante entre modelos: a teoria de equilíbrio \textbf{jamais} poderia
produzi-las, porque não tem não-linearidade alguma.

Um modelo dinâmico pode produzi-las — mas só se conservar as não-linearidades
que as geram. Vale registrar que a formulação de \S\ref{sec:aguasrasas} não as
conserva: a advecção foi descartada por $\mathrm{Ro}\ll 1$ e a continuidade foi
linearizada com $H$ em lugar de $H+\zeta$. Sobra apenas o atrito quadrático de
fundo, que gera sobremarés fracamente. Produzir M\textsubscript{4} e
M\textsubscript{6} com realismo exigiria \emph{reintroduzir} os termos
não-lineares \emph{e} resolver os estuários onde elas vivem — duas condições, e
não uma.

% =====================================================================
\chapter{Métodos numéricos}
\label{cap:numerica}

A formulação dos capítulos anteriores só vira resultado através de escolhas de
discretização, e várias delas têm conteúdo físico. Este capítulo trata das que
importam.

\section{Grade C de Arakawa}
\label{sec:gradec}

As variáveis não ficam todas no mesmo ponto. Na \textbf{grade C} de Arakawa
\cite{arakawa1977}, a elevação $\zeta$ mora no centro da célula, $u$ na face
leste e $v$ na face norte:

\begin{figure}[htbp]
\centering
\begin{tikzpicture}[scale=1.5]
  \draw[gray!50] (0,0) grid (3,2);
  \foreach \x in {0.5,1.5,2.5}{\foreach \y in {0.5,1.5}{
    \fill[azulforte] (\x,\y) circle (1.6pt);
    \node[font=\scriptsize,azulforte,above right=-1pt] at (\x,\y) {$\zeta$};}}
  \foreach \x in {0,1,2,3}{\foreach \y in {0.5,1.5}{
    \draw[laranja,thick,-Latex] (\x,\y) -- ++(0.3,0);
    \node[font=\scriptsize,laranja,below=1pt] at (\x,\y) {$u$};}}
  \foreach \x in {0.5,1.5,2.5}{\foreach \y in {0,1,2}{
    \draw[verde,thick,-Latex] (\x,\y) -- ++(0,0.3);
    \node[font=\scriptsize,verde,right=1pt] at (\x,\y) {$v$};}}
\end{tikzpicture}
\caption{Arranjo da grade C. O gradiente de pressão e a divergência do fluxo são
calculados com diferenças de \emph{um} intervalo entre pontos adjacentes, sem
média intermediária. Só o termo de Coriolis exige interpolação — de quatro
vizinhos.}
\label{fig:gradec}
\end{figure}

A justificativa é a representação das ondas de gravidade-inércia. A grade C
representa corretamente a dispersão \eqref{eq:poincare} quando
$\Delta x\lesssim L_R$, e é a melhor escolha nesse regime. Aqui
$\Delta x=\SI{110}{\kilo\meter}$ contra $L_R\approx\SI{1900}{\kilo\meter}$
(\S\ref{sec:ondas}): estamos com folga de mais de uma ordem de grandeza no
regime favorável.

O preço é que o termo de Coriolis precisa de $v$ nos pontos $u$ e vice-versa,
obtidos pela média de quatro vizinhos. A alternativa (grade B) inverte o
compromisso: Coriolis fica exato, o gradiente de pressão passa a exigir média, e
a dispersão degrada em $\Delta x>L_R$.

\section{Avanço forward-backward}
\label{sec:fb}

O esquema temporal é o \emph{forward-backward}: a continuidade avança primeiro,
usando as velocidades \textbf{antigas}; em seguida o momento avança usando a
elevação \textbf{recém-calculada}.
\begin{align}
\zeta^{n+1}&=\zeta^{n}-\Delta t\,\nabla\!\cdot\!(H\vect{u}^{n}),\\
u^{n+1}&=u^{n}+\Delta t\left[\;\;f\,v^{n}
  -\frac{g}{a\cos\phi}\pd{\zeta'^{\,n+1}}{\lambda}+\dots\right],\\
v^{n+1}&=v^{n}+\Delta t\left[-f\,u^{n+1}
  -\frac{g}{a}\pd{\zeta'^{\,n+1}}{\phi}+\dots\right].
\end{align}
Repare no $u^{n+1}$ dentro do Coriolis de $v$: o mesmo princípio de usar sempre
o valor mais recente disponível, aplicado uma segunda vez dentro do passo de
momento.

As propriedades que motivam a escolha:

\begin{itemize}
\item \textbf{estabilidade neutra} para ondas de gravidade: o fator de
  amplificação tem módulo exatamente 1 dentro do limite CFL, sem amortecimento
  numérico artificial. Isso importa num problema em que a dissipação física é
  justamente o que se quer calibrar (\S\ref{sec:calibracao});
\item \textbf{segunda ordem efetiva} em $\Delta t$ para o modo de gravidade,
  ao custo de \emph{uma} avaliação por passo;
\item \textbf{sem modo computacional}: ao contrário do leapfrog, não há solução
  espúria alternante que exija filtro de Asselin — filtro que, por sua vez,
  introduziria amortecimento espúrio.
\end{itemize}

\section{Atrito quadrático semi-implícito}

O termo $C_d|\vect{u}|u/H$ tratado explicitamente instabiliza-se quando
$\Delta t\,C_d|\vect{u}|/H>1$, o que ocorre em células rasas com corrente forte
— exatamente onde a dissipação é maior. A solução usada é semi-implícita:
avalia-se o módulo $|\vect{u}|$ no tempo antigo e a velocidade linear no novo,
\begin{equation}
u^{n+1}=\frac{u^{n}+\Delta t\,\mathcal{R}}
{1+\Delta t\left(\dfrac{C_d|\vect{u}^{n}|}{H}+r\right)},
\label{eq:semiimplicito}
\end{equation}
com $\mathcal{R}$ os demais termos. O denominador é sempre $>1$, de modo que o
esquema é \textbf{incondicionalmente dissipativo e limitado}: nenhum valor de
$\Delta t$ pode fazer a velocidade trocar de sinal por excesso de atrito, que é
o modo de falha do tratamento explícito.

\section{Condição CFL e o limite polar}
\label{sec:cfl}

Para ondas de gravidade em duas dimensões, a condição de Courant é
\begin{equation}
\Delta t\;\le\;\frac{\min(\Delta x,\Delta y)}{\sqrt{2\,gH}} .
\end{equation}
Adota-se $\Delta t=s\cdot\min(\Delta x,\Delta y)/\sqrt{gH}$ com fator de
segurança $s\in[\num{0.40},\num{0.45}]$, abaixo do limite estrito
$1/\sqrt2=\num{0.707}$.

Como $\Delta x=a\Delta\lambda\cos\phi$, a célula encolhe com a latitude e o
passo estável junto. Este é o motivo do corte em \ang{86}
(\S\ref{sec:batimetria}): na última linha de células molhadas
$\Delta x\approx\SI{8.7}{\kilo\meter}$ e
$\Delta t\approx\SI{19}{\second}$, contra \SI{110}{\kilo\meter} e
$\sim\SI{250}{\second}$ no equador. \textbf{Uma única célula polar governa o
custo da rodada inteira} — a observação prática que justifica a decisão de
projeto.

\begin{observacao}[O caminho não tomado]
Modelos globais de alta resolução resolvem o problema polar de outras formas:
grade tripolar (com os polos numéricos deslocados para cima de continentes),
malha não estruturada, ou filtragem zonal em alta latitude. Todas são
preferíveis; nenhuma é barata. Cortar em \ang{86} é a solução aceitável para um
modelo desta escala, \emph{desde que declarada}.
\end{observacao}

\section{Integradores: a pergunta sobre RK4, respondida por medida}
\label{sec:integradores}

Uma pergunta que se impõe ao montar o modelo é se Runge--Kutta de quarta
ordem dá conta, ou se outro método seria mais adequado. A resposta correta é
que \textbf{há dois subsistemas com exigências opostas}, e ela difere para
cada um.

\subsection{Resposta da bacia: RK4 é a escolha certa}

A Eq.~\eqref{eq:oscilador} tem amortecimento $2\gamma\dot h$ que
\textbf{depende da velocidade}. O sistema deixa de ser hamiltoniano, de modo que
a análise de \S\ref{sec:simpletica} não se aplica: não há estrutura simplética a
preservar, e os métodos que a preservam perdem sua vantagem — além de deixarem
de ser explícitos, já que a força passa a depender da incógnita que estão
avançando. O sistema não é rígido
($\Delta t=\SI{60}{\second}$ contra período de \SI{12.42}{\hour} dá
\num{745} passos por ciclo) e há forçamento externo dependente do tempo, que o
RK4 avalia em $t+\Delta t/2$. É o caso de livro para RK4.

\begin{alerta}[A armadilha do forçamento pré-amostrado]
A quarta ordem do RK4 vem justamente das avaliações em $t+\Delta t/2$. Se
$h_{\text{eq}}$ for passada como \emph{array} amostrado apenas nos instantes
$t$, em vez de como função, o método cai \textbf{silenciosamente} para segunda
ordem: nada falha, nada avisa, e o erro simplesmente fica maior do que deveria.
A implementação precisa, portanto, receber o forçamento como \emph{função do
tempo}, e não como vetor pré-amostrado — e a ordem observada deve ser medida, e
não suposta.
\end{alerta}

\subsection{Órbitas: RK4 é o pior dos quatro}

Aqui vale \S\ref{sec:simpletica}: o RK4 não é simplético, seu erro de energia
tem componente secular, e a Eq.~\eqref{eq:amplificacao3x} amplifica por três
qualquer desvio em $D$.

\begin{table}[htbp]
\centering
\caption{Órbita lunar integrada por \SI{10}{\anos}, contra o Kepler analítico.
``Erro em $h$'' aplica
\eqref{eq:amplificacao3x} ao erro de distância.}
\label{tab:integradores}
\footnotesize
\begin{tabular}{@{}lrrrrrrl@{}}
\toprule
método & $\Delta t$ [s] & $|\delta E/E|$ & erro pos.\ [m] & erro em $h$ [\%]
& aval./passo & tempo [s] & secular? \\
\midrule
RK4             & 600  & \num{3.65e-13} & \num{1.37e-1} & \num{1.26e-9} & 4 & \num{9.6} & \textbf{sim} \\
Velocity-Verlet & 600  & \num{1.15e-7}  & \num{3.22e5}  & \num{4.87e-3} & 1 & \num{3.4} & não \\
Yoshida-4       & 600  & \num{3.86e-13} & \num{2.00e0}  & \num{1.22e-8} & 3 & \num{7.7} & não \\
\addlinespace
RK4             & 1800 & \num{6.64e-11} & \num{2.05e1}  & \num{4.20e-7} & 4 & \num{3.2} & \textbf{sim} \\
Velocity-Verlet & 1800 & \num{1.03e-6}  & \num{2.90e6}  & \num{4.38e-2} & 1 & \num{1.1} & não \\
Yoshida-4       & 1800 & \num{3.42e-11} & \num{1.64e2}  & \num{8.87e-7} & 3 & \num{2.5} & não \\
\addlinespace
RK4             & 3600 & \num{2.10e-9}  & \num{5.86e2}  & \num{1.77e-5} & 4 & \num{1.6} & \textbf{sim} \\
Velocity-Verlet & 3600 & \num{4.15e-6}  & \num{1.16e7}  & \num{1.75e-1} & 1 & \num{0.6} & não \\
Yoshida-4       & 3600 & \num{5.49e-10} & \num{2.62e3}  & \num{1.42e-5} & 3 & \num{1.3} & não \\
\bottomrule
\end{tabular}
\end{table}

A ordem de convergência observada confirma a declarada: RK4 \num{4.023},
Velocity-Verlet \num{2.000}, Yoshida-4 \num{3.998}. E a medida que separa
secular de limitado — a razão entre o erro de energia em \SI{10}{\anos} e em
\SI{1}{\ano}, a $\Delta t=\SI{3600}{\second}$ — dá \num{8.9} para o RK4 contra
\num{0.9} para ambos os simpléticos. Dito de outro modo: o RK4 passa
\SI{39}{\percent} do tempo estabelecendo um novo máximo de erro; os
simpléticos, \SI{3}{\percent}.

\subsection{Yoshida-4: quarta ordem \emph{e} erro limitado}

A composição de Yoshida \cite{yoshida1990} constrói um método simplético de
quarta ordem a partir de três saltos de Verlet com passos
$w_1\Delta t,\;w_0\Delta t,\;w_1\Delta t$. Impondo consistência
($2w_1+w_0=1$) e cancelamento do termo de erro de terceira ordem
($2w_1^{3}+w_0^{3}=0$):
\begin{equation}
w_1=\frac{1}{2-2^{1/3}}=\num{1.35120719},\qquad
w_0=\frac{-2^{1/3}}{2-2^{1/3}}=\num{-1.70241438}.
\end{equation}
Note que $w_0<0$: o método \textbf{anda para trás no tempo} no salto central.
Isso é inevitável — não existe composição simplética de ordem $\ge 3$ com todos
os passos positivos — e é o preço de combinar ordem alta com simpleticidade.

Yoshida-4 vence: quarta ordem, erro limitado, e três avaliações de força por
passo contra quatro do RK4. É simultaneamente mais estável e mais barato.

\subsection{A escolha adotada, e o veredito}
\label{sec:kepleranalitico}

Adota-se o \textbf{Kepler analítico} (\S\ref{cap:gravitacao}): precisão de
máquina, deriva exatamente zero, custo $\ordem{1}$ por instante consultado — e,
decisivamente, permite avaliar a posição em um $t$ \emph{arbitrário}, o que é
necessário porque o RK4 da bacia pede a efeméride em $t+\Delta t/2$. Os demais
integradores foram implementados e mantidos disponíveis, porque é isso que torna
a Tabela~\ref{tab:integradores} reproduzível.

\begin{caixa}[O veredito, com a magnitude na frente]
O pior erro de maré de toda a Tabela~\ref{tab:integradores} é
\SI{0.175}{\percent}, ou seja \textbf{\SI{0.95}{\milli\meter} numa maré de
\SI{0.54}{\meter}} — e vem do método mais fraco no passo mais grosseiro. Para
esta aplicação, portanto, \textbf{a escolha do integrador é numericamente
livre}: a sugestão original de usar RK4 funcionaria perfeitamente, e nenhum
resultado físico mudaria.

Isso não invalida a análise; delimita-a. A distinção
simplético/não-simplético é real e mensurável, e seria decisiva num problema
orbital de escala maior (integração de sistemas planetários por $10^{8}$ anos,
por exemplo). Saber \emph{quando} um argumento correto é irrelevante é parte do
que se espera de uma análise numérica bem feita.
\end{caixa}

\section{Métodos considerados e descartados}

\begin{description}[style=nextline,leftmargin=0.8cm]
\item[Runge--Kutta encaixados adaptativos (RKF45, DOPRI5, DOP853)] Usados apenas como referência de alta precisão
  ($\text{rtol}=10^{-12}$). O
  passo adaptativo só compensaria se $D$ variasse muito dentro de um passo, o
  que não ocorre com $e=\num{0.055}$.
\item[Leapfrog puro] Equivalente ao Velocity-Verlet em precisão, mas com
  velocidades meio passo fora de fase, o que atrapalharia o cálculo de energia e
  a interpolação de efemérides.
\item[Métodos implícitos (Radau, BDF)] Feitos para sistemas rígidos. Nenhum dos
  subsistemas aqui é rígido: as escalas de tempo relevantes vão de \SI{12}{\hour}
  a \SI{29.5}{\day}, uma separação de apenas $\sim\!60\times$.
\end{description}

% =====================================================================
\part{Validação, limites e síntese}
\label{parte:sintese}
% =====================================================================

\chapter{Validação}
\label{cap:validacao}

Um modelo geofísico só existe como afirmação verificável se for confrontado com
observação independente. Este capítulo descreve as duas confrontações feitas, e
— tão importante quanto — o critério de escolha dos alvos.

\section{Por que estações insulares}
\label{sec:porqueilhas}

Numa grade de \ang{1.0}, uma célula tem \SI{110}{\kilo\meter}. Portos,
estuários e plataformas estreitas simplesmente \textbf{não existem} nessa
grade. Comparar o modelo com um marégrafo de estuário mediria a resolução da
grade, não a física do modelo.

Ilhas oceânicas isoladas, ao contrário, amostram a maré de oceano aberto —
precisamente o que um modelo barotrópico grosseiro tem chance de acertar. São
elas o alvo justo.

\begin{caixa}[Metodologia: escolher o alvo antes de ver o resultado]
A escolha das oito ilhas foi feita \emph{a priori}, por um critério físico
(amostrar oceano aberto), e não selecionando as estações em que o modelo se
saía bem. Além disso, cinco estações costeiras foram incluídas
\textbf{deliberadamente para mostrar onde o modelo falha} — não para inflar a
estatística. As duas populações são reportadas separadamente, com RMS
separados. Um resultado agregado sobre as treze estações seria menos
informativo e mais lisonjeiro, e por isso não é o que se reporta.
\end{caixa}

\section{Amplitude de M\textsubscript{2} contra a NOAA CO-OPS}

\begin{table}[htbp]
\centering
\caption{Modelo a \ang{1.0} contra constantes harmônicas publicadas pela
NOAA CO-OPS.}
\label{tab:validacao}
\begin{tabular}{@{}llrrrr@{}}
\toprule
estação & tipo & obs.\ [m] & modelo [m] & erro [m] & erro [\%] \\
\midrule
Hilo, Hawaii        & ilha & \num{0.220} & \num{0.219} & \num{-0.001} & \num{-0.4} \\
Apra Harbor, Guam   & ilha & \num{0.220} & \num{0.216} & \num{-0.004} & \num{-1.6} \\
Sand Island, Midway & ilha & \num{0.119} & \num{0.111} & \num{-0.008} & \num{-6.5} \\
Nawiliwili, Kauai   & ilha & \num{0.149} & \num{0.161} & \num{0.012}  & \num{8.3} \\
Kwajalein           & ilha & \num{0.467} & \num{0.422} & \num{-0.045} & \num{-9.5} \\
Honolulu, HI        & ilha & \num{0.171} & \num{0.149} & \num{-0.022} & \num{-13.0} \\
Wake Island         & ilha & \num{0.277} & \num{0.225} & \num{-0.052} & \num{-18.7} \\
Pago Pago, Samoa    & ilha & \num{0.377} & \num{0.506} & \num{0.129}  & \num{34.2} \\
\midrule
The Battery, NY     & costeira & \num{0.671} & \num{0.658} & \num{-0.013} & \num{-1.9} \\
Key West, FL        & costeira & \num{0.179} & \num{0.162} & \num{-0.017} & \num{-9.3} \\
Boston, MA          & costeira & \num{1.371} & \num{0.966} & \num{-0.405} & \num{-29.5} \\
San Francisco, CA   & costeira & \num{0.576} & \num{0.392} & \num{-0.184} & \num{-32.0} \\
Anchorage, AK       & costeira & \num{3.505} & \num{0.878} & \num{-2.627} & \num{-75.0} \\
\midrule
\multicolumn{2}{@{}l}{\textbf{RMS ilhas}} & \multicolumn{4}{l}{\textbf{\SI{5.2}{\centi\meter}} sobre 8 estações (obs.\ média \SI{25.0}{\centi\meter})} \\
\multicolumn{2}{@{}l}{\textbf{RMS costeiras}} & \multicolumn{4}{l}{\SI{119.2}{\centi\meter} sobre 5 estações (obs.\ média \SI{126.0}{\centi\meter})} \\
\bottomrule
\end{tabular}
\end{table}

\begin{figure}[htbp]
\centering
\includegraphics[width=0.86\textwidth]{validacao_M2.png}
\caption{Modelo contra observação em escala logarítmica, com guias de fator 2.
A escala log-log é a escolha correta aqui porque as amplitudes cobrem duas
ordens de grandeza; num eixo linear as estações insulares — o alvo justo do
modelo — ficariam todas comprimidas na origem.}
\label{fig:validacao}
\end{figure}

\subsection{Taxonomia dos erros}

Os erros da Tabela~\ref{tab:validacao} não são aleatórios; separam-se em
categorias com causas identificáveis:

\begin{description}[style=nextline,leftmargin=0.8cm]
\item[Erro de resolução (Anchorage, $-\SI{75}{\percent}$)] O Cook Inlet tem
  \SI{30}{\kilo\meter} de largura. Numa grade de \SI{110}{\kilo\meter} ele não
  existe, e o mecanismo de amplificação que produz \SI{3.5}{\meter} está
  fisicamente ausente. Aumentar a resolução resolveria.
\item[Erro de plataforma (Boston, San Francisco, $\sim\!-\SI{30}{\percent}$)]
  A plataforma continental é resolvida de forma grosseira; a lei de Green
  ($A\propto H^{-1/4}$) é aplicada a uma profundidade média que suaviza o
  gradiente real.
\item[Erro aceitável (The Battery, Key West, $<\SI{10}{\percent}$)]
  Curiosamente boas, porque estão em plataformas largas o bastante para a grade
  enxergar.
\item[Erro de resposta local (Pago Pago, $+\SI{34}{\percent}$)] Único erro
  positivo grande entre as ilhas. Pago Pago fica numa baía estreita cuja
  ressonância local o modelo não pode representar — aqui o modelo erra
  \emph{para cima} porque a célula amostra oceano aberto adjacente numa região
  de alta amplitude.
\end{description}

\section{A máquina harmônica contra as tábuas oficiais}

Independentemente do modelo dinâmico, o módulo de análise harmônica foi
verificado contra as \textbf{previsões oficiais da NOAA}: dadas as constantes
que a própria NOAA publica, reconstruímos as tábuas de maré deles com
\eqref{eq:previsao}.

\begin{center}
\begin{tabular}{@{}lrrr@{}}
\toprule
estação & amplitude observada & RMS da diferença & correlação \\
\midrule
Honolulu & \SI{0.72}{\meter} & \SI{0.052}{\meter} & \num{0.991} \\
San Francisco & \SI{2.11}{\meter} & \SI{0.067}{\meter} & \num{0.994} \\
Boston & \SI{3.69}{\meter} & \SI{0.110}{\meter} & \num{0.995} \\
Key West & \SI{0.70}{\meter} & \SI{0.078}{\meter} & \num{0.990} \\
\bottomrule
\end{tabular}
\end{center}

E no sentido inverso — ajustando as previsões da NOAA com os nossos argumentos
astronômicos — recuperam-se as constantes publicadas com amplitude exata a três
casas decimais e fase a menos de \ang{1}:

\begin{center}
\begin{tabular}{@{}lcccccccc@{}}
\toprule
& M\textsubscript{2} & S\textsubscript{2} & N\textsubscript{2} & K\textsubscript{2}
& K\textsubscript{1} & O\textsubscript{1} & P\textsubscript{1} & Q\textsubscript{1} \\
\midrule
$\Delta g$ & $+\ang{0.0}$ & $-\ang{0.2}$ & $+\ang{0.8}$ & $+\ang{0.4}$
& $-\ang{0.1}$ & $+\ang{0.1}$ & $+\ang{0.4}$ & $+\ang{0.8}$ \\
\bottomrule
\end{tabular}
\end{center}

O resíduo de 5 a \SI{11}{\centi\meter} vem das $\sim\!29$ constituintes que a
NOAA usa e o modelo não inclui. Foi este teste que revelou o erro de
\ang{180} nas diurnas (\S\ref{sec:correcaofase}).

\begin{caixa}[Por que esta validação é independente]
Ela não usa o modelo dinâmico em nenhum ponto: testa apenas a máquina de
argumentos astronômicos, fatores nodais e convenção de fase. Separar as duas
validações significa que, se a comparação da Tabela~\ref{tab:validacao} tivesse
saído ruim, saberíamos que a culpa seria da hidrodinâmica e não da contabilidade
de fases — e vice-versa. Testar subsistemas contra fontes externas distintas é
o que permite localizar a falha em vez de apenas detectá-la.
\end{caixa}

\section{O que emerge sem ser programado}

Nenhuma frequência de maré, nenhum ponto anfidrômico e nenhum sentido de
rotação está codificado. Tudo sai das constantes orbitais e das equações:

\begin{itemize}
\item \textbf{M\textsubscript{2} em \SI{12}{\hour}\,\SI{25.2}{\minute}},
  S\textsubscript{2} em \SI{12}{\hour} exatas, batimento de sizígia/quadratura
  em \SI{14.77}{\day}, desigualdade diurna — todos verificados por FFT da série
  simulada;
\item \textbf{pontos anfidrômicos} nos lugares certos, detectados pelo critério
  topológico \eqref{eq:winding};
\item \textbf{o sentido de rotação segue Coriolis} (Tabela~\ref{tab:rotacao});
\item \textbf{amplificação em água rasa} — as maiores amplitudes aparecem em
  plataforma, não em oceano profundo, conforme a lei de Green;
\item \textbf{o ciclo nodal de \SI{18.6}{\anos}} — com o nó precessando, as
  diurnas modulam $\pm\SI{18}{\percent}$, o que a análise harmônica corrige com
  os fatores nodais.
\end{itemize}

% =====================================================================
\chapter{Limites e a escada até os modelos operacionais}
\label{cap:limites}

\section{Onde este modelo está}

\begin{table}[htbp]
\centering
\caption{Níveis de sofisticação em modelagem de maré e a acurácia
correspondente em oceano profundo \cite{stammer2014}.}
\label{tab:escada}
\small
\begin{tabular}{@{}p{4.3cm}p{5.6cm}r@{}}
\toprule
nível & o que acrescenta & RMS \\
\midrule
0 --- equilíbrio (Newton) & potencial de grau 2, resposta estática & forma funcional errada \\
1 --- \textbf{LTE barotrópicas} $\leftarrow$ \emph{aqui} & Coriolis, batimetria, atrito, SAL & dezenas de cm \\
2 --- LTE de alta resolução & \ang{0.0625}, malha não estruturada & $\sim\SI{10}{\centi\meter}$ \\
3 --- assimilação de altimetria & décadas de TOPEX/Jason & $\sim\SI{1}{\centi\meter}$ \\
\bottomrule
\end{tabular}
\end{table}

Os modelos operacionais (FES2014/2022, TPXO9/10, GOT5, EOT20, DTU23) estão
todos no nível 3 e reportam $\approx\SI{0.9}{\centi\meter}$ de RMS em oceano
profundo, $\approx\SI{5}{\centi\meter}$ em plataforma continental e
$\approx\SI{6.5}{\centi\meter}$ em água costeira \cite{stammer2014,lyard2021}.

\section{Termos de física ainda ausentes}

\begin{description}[style=nextline,leftmargin=0.8cm]
\item[Maré interna resolvida] Aqui ela é apenas um arrasto linear calibrado
  \eqref{eq:arrastolinear}. Resolvê-la exige estratificação e um modelo
  baroclínico — outro modelo, não um refinamento deste.
\item[SAL completo] A aproximação escalar $\beta=\num{0.09}$ erra até
  \SI{20}{\percent} em plataforma (\S\ref{cap:sal}); o cálculo correto é a
  convolução global \eqref{eq:sal_exato} a cada passo.
\item[Maré atmosférica e barômetro invertido] $\sim\SI{1}{\centi\meter}$ por
  hPa de pressão.
\item[Gelo marinho] no Ártico e no Antártico, que modifica o atrito superficial.
\item[Termos não-lineares] a advecção $(\vect{u}\cdot\nabla)\vect{u}$ e o
  $(H+\zeta)$ da continuidade, ambos descartados na linearização
  (\S\ref{sec:escala}). São eles que geram M\textsubscript{4} e
  M\textsubscript{6} em estuários; com a formulação linearizada, só o atrito
  quadrático de fundo pode gerá-las, e fracamente.
\end{description}

Nada disso é obstáculo conceitual: é escala computacional e dados. O salto
qualitativo — o que separa ``bojo que segue a Lua'' de ``onda que gira em torno
de anfidromos'' — já está feito.

\section{Síntese: a cadeia completa do argumento}

Vale reconstruir, em uma página, o encadeamento que o documento percorreu.

\begin{enumerate}
\item O referencial geocêntrico é não inercial; o termo de arrasto
  $-m\ddot{\vect{R}}_0$ cancela exatamente o termo $n=1$ da expansão do campo
  perturbador (Cap.~\ref{cap:referenciais}--\ref{cap:gravitacao}).
\item Sobra o \textbf{quadrupolo}, $\propto D^{-3}$. Daí a Lua vencer o Sol
  apesar de ser 178 vezes mais fraca em atração direta
  (Cap.~\ref{cap:forcamare}--\ref{cap:potencial}).
\item Supondo resposta instantânea, a superfície é equipotencial e
  $h=\SI{0.54}{\meter}$ com dois bojos antipodais — a teoria de equilíbrio, que
  reproduz todos os números do texto-base (Cap.~\ref{cap:potencial}).
\item Mas a onda de gravidade viaja a \SI{198}{\meter\per\second} contra os
  \SI{448}{\meter\per\second} do ponto sublunar: \textbf{a resposta não pode ser
  instantânea} (Cap.~\ref{cap:falha}).
\item Portanto o problema correto é hidrodinâmico: águas rasas em esfera
  girante, com batimetria, Coriolis, dissipação e auto-atração — as Equações de
  Maré de Laplace (Cap.~\ref{cap:lte}--\ref{cap:dissipacao}).
\item Sob Coriolis, a onda forçada e refletida organiza-se em torno de
  \textbf{defeitos topológicos do campo de fase} — os pontos anfidrômicos —,
  cuja existência a teoria de equilíbrio é estruturalmente incapaz de prever
  (Cap.~\ref{cap:anfidromos}).
\item A dissipação desalinha o bojo, o desalinhamento gera \textbf{torque}, e o
  torque realimenta a mecânica celeste do item 1: a Lua se afasta, o dia se
  alonga (\S\ref{sec:torque}).
\end{enumerate}

O círculo se fecha. A maré começa como um efeito de referencial na dinâmica de
dois corpos e termina alterando esse mesmo problema de dois corpos.

% =====================================================================
\appendix
% =====================================================================

\chapter{Teorema da adição e as três espécies}
\label{ap:legendre}

A fatoração \eqref{eq:especies}, que separa geografia de astronomia no
potencial de maré, decorre do teorema da adição dos harmônicos esféricos:
\begin{equation}
P_n(\cos\psi)=P_n(\cos\theta_1)P_n(\cos\theta_2)
+2\sum_{m=1}^{n}\frac{(n-m)!}{(n+m)!}
P_n^{m}(\cos\theta_1)P_n^{m}(\cos\theta_2)\cos\!\big(m(\varphi_1-\varphi_2)\big),
\end{equation}
com $\psi$ o ângulo entre as direções $(\theta_1,\varphi_1)$ e
$(\theta_2,\varphi_2)$.

Identificando a primeira direção com o observador (colatitude
$\theta_1=\pi/2-\phi$, de modo que $\cos\theta_1=\sin\phi$) e a segunda com o
corpo perturbador ($\cos\theta_2=\sin\delta$), e chamando de $H$ a diferença de
longitudes — o ângulo horário $\mathcal{H}$ —, o caso $n=2$ dá coeficientes
\begin{equation}
m=1:\;2\cdot\frac{1!}{3!}=\frac13,
\qquad
m=2:\;2\cdot\frac{0!}{4!}=\frac1{12},
\end{equation}
que são exatamente os fatores de \eqref{eq:especies}. Com
\begin{equation}
P_2(x)=\tfrac12(3x^{2}-1),\qquad
P_2^{1}(x)=3x\sqrt{1-x^{2}},\qquad
P_2^{2}(x)=3(1-x^{2}),
\end{equation}
e $x=\sin\phi$, as dependências latitudinais tornam-se
\begin{align}
m=0:&\quad \tfrac12(3\sin^{2}\phi-1) &&\text{(longo período)}\\
m=1:&\quad 3\sin\phi\cos\phi=\tfrac32\sin 2\phi &&\text{(diurna)}\\
m=2:&\quad 3\cos^{2}\phi &&\text{(semidiurna)}
\end{align}
— as três formas usadas ao longo do texto, a menos de constantes absorvidas
na amplitude tabelada de cada constituinte.

\begin{observacao}
As dependências em \emph{declinação} são as mesmas funções aplicadas a
$\sin\delta$, e é daí que saem as regras práticas: a espécie diurna
$\propto\sin 2\delta$ some quando a Lua cruza o equador; a semidiurna
$\propto\cos^{2}\delta$ é máxima então. Um mesmo teorema explica a estrutura
espacial e a modulação temporal.
\end{observacao}

% ---------------------------------------------------------------------
\chapter{Derivação das equações de águas rasas}
\label{ap:sw}

Parta da incompressibilidade $\nabla\cdot\vect{u}_3=0$ e integre em $z$ de
$-H(\vect{x})$ a $\zeta(\vect{x},t)$:
\begin{equation}
\int_{-H}^{\zeta}\left(\pd{u}{x}+\pd{v}{y}\right)\dd z
+ w\big|_{\zeta}-w\big|_{-H}=0 .
\end{equation}

As condições cinemáticas são
\begin{equation}
w\big|_{\zeta}=\pd{\zeta}{t}+\vect{u}\cdot\nabla\zeta,
\qquad
w\big|_{-H}=-\vect{u}\cdot\nabla H
\end{equation}
(a superfície é material; o fundo é impermeável). Usando a regra de Leibniz
para trocar derivada e integral, e a independência de $\vect{u}$ em relação a
$z$ — que é consequência do equilíbrio hidrostático com $\rho$ constante
(\S\ref{sec:aguasrasas}) —, os termos de fronteira cancelam e resta
\begin{equation}
\pd{\zeta}{t}+\nabla\cdot\big[(H+\zeta)\vect{u}\big]=0 .
\end{equation}

A linearização $H+\zeta\to H$ vale porque $\zeta/H\sim\num{2.5e-4}$ em oceano
aberto; ela deixa de valer em água rasa, e é precisamente essa não-linearidade
que gera as constituintes de águas rasas.

Em coordenadas esféricas, com os fatores métricos
$\dd x=a\cos\phi\,\dd\lambda$ e $\dd y=a\,\dd\phi$, o divergente adquire o
fator $1/(a\cos\phi)$ e o $\cos\phi$ interno da Eq.~\eqref{eq:lte_cont}:
\begin{equation}
\nabla\cdot\vect{A}=\frac{1}{a\cos\phi}
\left[\pd{A_\lambda}{\lambda}+\pd{(A_\phi\cos\phi)}{\phi}\right].
\end{equation}

\begin{observacao}[O detalhe discreto que corresponde a isso]
Na grade, o $\cos\phi$ interno é avaliado na \emph{face norte} da célula
($\phi+\Delta\phi/2$), e não no centro. Avaliá-lo no centro quebraria a
conservação de volume: o fluxo que sai de uma célula deixaria de ser
exatamente o que entra na vizinha. É o tipo de detalhe que não altera a
convergência formal mas destrói a conservação.
\end{observacao}

% ---------------------------------------------------------------------
\chapter{Implementação em Python}
\label{ap:python}

Os resultados numéricos deste texto foram produzidos por uma implementação em
Python (pacote \texttt{oceantides}), com NumPy para as operações de grade,
SciPy apenas como referência de alta precisão nos testes e Matplotlib para as
figuras. Este apêndice resume como o modelo está organizado e mostra os três
trechos em que a tradução da equação para o código não é imediata.

\section*{Organização}

\begin{table}[H]
\centering
\small
\begin{tabular}{@{}lll@{}}
\toprule
módulo & conteúdo & seção correspondente \\
\midrule
\texttt{constants.py}   & constantes físicas e elementos orbitais & --- \\
\texttt{forcing.py}     & força de maré exata e expandida, potencial, $h_{\text{eq}}$ & Cap.~\ref{cap:forcamare}--\ref{cap:potencial} \\
\texttt{ephemeris.py}   & posição da Lua e do Sol; equação de Kepler & \S\ref{sec:kepler_eq} \\
\texttt{integrators.py} & RK4, Velocity-Verlet, Yoshida-4 & \S\ref{sec:integradores} \\
\texttt{response.py}    & oscilador forçado de bacia & Cap.~\ref{cap:oscilador} \\
\texttt{bathymetry.py}  & batimetria ETOPO, máscara de terra, limite CFL & \S\ref{sec:batimetria} \\
\texttt{lte.py}         & Equações de Maré de Laplace, grade C & Cap.~\ref{cap:lte} \\
\texttt{harmonic.py}    & constituintes, fatores nodais, anfidromos & Cap.~\ref{cap:harmonica} \\
\texttt{validation.py}  & comparação com marégrafos & Cap.~\ref{cap:validacao} \\
\bottomrule
\end{tabular}
\end{table}

O campo de maré nunca é armazenado: a maré de equilíbrio é forma fechada,
recalculada quando pedida, e a solução das Equações de Maré de Laplace guarda
apenas os coeficientes harmônicos, dos quais o campo em qualquer instante é
reconstruído. Em vez de girar os $\sim\!\num{65000}$ pontos da grade para o
referencial inercial a cada passo, gira-se o vetor que aponta para o corpo
perturbador — um vetor em vez de uma grade.

\section*{O passo de tempo}

O esquema \emph{forward-backward} de \S\ref{sec:fb} e o atrito semi-implícito
de \eqref{eq:semiimplicito}, na forma em que são executados:

\begin{lstlisting}[numbers=none]
# continuidade: eta avanca com as velocidades antigas
div = (flux_u - np.roll(flux_u, 1, axis=1)) / (R * cos_c * dlambda)
eta = np.where(mask, eta - dt * div, 0.0)

# desvio de equipotencial: SAL + forcante astronomica
eta_dev = (1.0 - sal_beta) * eta - self.equilibrium_field(t + dt)

# momento zonal, ja com o eta novo; atrito no denominador
rhs_u = u + dt * (f_c * v_at_u - g * dpdx)
u = rhs_u / (1.0 + dt * (drag * speed_u / h_u + internal_drag))
\end{lstlisting}

\section*{Análise harmônica acumulada}

A matriz normal \eqref{eq:normal} é construída durante a integração, sem que a
série temporal seja guardada (\S\ref{sec:embutida}). O custo de memória é de
$\ordem{n_{\text{con}}}$ campos, e não $\ordem{n_{\text{passos}}}$:

\begin{lstlisting}[numbers=none]
basis = np.empty(2 * n_con)
basis[0::2] = np.cos(sigma * ts)
basis[1::2] = np.sin(sigma * ts)
acc    += basis[:, None, None] * eta           # elevacao
acc_u  += basis[:, None, None] * uc            # corrente zonal
acc_v  += basis[:, None, None] * vc            # corrente meridional
normal += np.outer(basis, basis)
\end{lstlisting}

\section*{Detecção de anfidromos}

O índice topológico \eqref{eq:winding} na grade discreta. A redução ao
intervalo principal, na quarta linha, é o que impede que a descontinuidade
$\ang{359}\to\ang{0}$ seja lida como uma volta:

\begin{lstlisting}[numbers=none]
winding = np.zeros_like(phase)
for k in range(len(ring)):
    diff = shifted[(k + 1) % len(ring)] - shifted[k]
    winding += ((diff + 180.0) % 360.0) - 180.0

is_amph = ocean & (np.abs(np.abs(winding) - 360.0) < 90.0) \
                & (amp <= amplitude_threshold)
\end{lstlisting}

\section*{Reprodução dos resultados}

\begin{lstlisting}[language=bash,numbers=none]
uv sync --extra dev --extra data
uv run --extra data python scripts/render_dynamics.py   # Cap. 12 e 15
uv run --extra data python scripts/tune_friction.py     # Tabela 11.1
uv run oceantides bench                                 # Tabela 14.1
uv run --extra data pytest                              # suite de testes
\end{lstlisting}

Cada resultado numérico citado no texto tem um teste correspondente que o trava:
os números do Capítulo~\ref{cap:potencial} contra os valores clássicos, a
resposta integrada do Capítulo~\ref{cap:oscilador} contra a admitância
analítica \eqref{eq:admitancia}, os períodos do Capítulo~\ref{cap:harmonica}
por FFT da série simulada, e as constantes harmônicas do
Capítulo~\ref{cap:validacao} contra a NOAA CO-OPS.

\chapter{Notação}
\label{ap:simbolos}

\begin{table}[H]
\centering
\small
\begin{tabular}{@{}lll@{}}
\toprule
símbolo & significado & valor típico \\
\midrule
$a$, $r_\oplus$ & raio médio da Terra & \SI{6.371e6}{\meter} \\
$g$ & gravidade padrão & \SI{9.80665}{\meter\per\second\squared} \\
$\Omega$ & velocidade angular da Terra & \SI{7.292115e-5}{\per\second} \\
$f=2\Omega\sin\phi$ & parâmetro de Coriolis & \SI{1.03e-4}{\per\second} a \ang{45} \\
$D$ & distância Terra--corpo & \SI{3.844e8}{\meter} (Lua) \\
$\psi$ & ângulo parcela--corpo, visto do centro & --- \\
$\phi,\lambda$ & latitude, longitude & --- \\
$\delta$, $\mathcal{H}$ & declinação e ângulo horário do corpo & --- \\
$H_i$, $g_i$ & amplitude e fase de Greenwich da constituinte $i$ & --- \\
$h$ (Doodson) & longitude média do Sol, Tab.~\ref{tab:doodson} & --- \\
$h_2$, $k_2$ & números de Love de maré & \num{0.61}, \num{0.30} \\
$\zeta$ & elevação da superfície do mar & \si{\meter} \\
$\zeta_{\text{eq}}$ & maré de equilíbrio (forçante) & \si{\meter} \\
$\zeta'$ & desvio de equipotencial, Eq.~\eqref{eq:desvio} & \si{\meter} \\
$H(\lambda,\phi)$ & profundidade do oceano & \SI{4000}{\meter} típica \\
$\vect{u}=(u,v)$ & velocidade média na vertical & \si{\meter\per\second} \\
$c=\sqrt{gH}$ & velocidade de fase da onda longa & \SI{198}{\meter\per\second} \\
$L_R=c/f$ & raio de deformação de Rossby & \SI{1900}{\kilo\meter} \\
$\sigma_c$ & velocidade angular da constituinte & \si{\radian\per\second} \\
$A$, $G$ & amplitude e fase de Greenwich & --- \\
$f_i$, $u_i$ & fatores nodais (amplitude e fase) & --- \\
$\beta$ & coeficiente escalar de SAL & \num{0.09} \\
$C_d$ & arrasto de fundo & \num{0.0025} \\
$r$ & arrasto linear (maré interna) & \SI{1e-5}{\per\second} \\
$k_2,h_2$ & números de Love de maré & \num{0.30}, \num{0.61} \\
$k'_n,h'_n$ & números de Love de carga & negativos \\
$\gamma_2=1+k_2-h_2$ & fator diminutivo & \num{0.69} \\
$Q=\omega_0/2\gamma$ & fator de qualidade da bacia & 2 a 20 \\
$w$ & número de voltas da fase, Eq.~\eqref{eq:winding} & $\pm 1$ \\
\bottomrule
\end{tabular}
\end{table}

\noindent
Três símbolos são reaproveitados, seguindo o uso corrente de cada literatura, e
convém tê-los presentes: $r$ é a distância radial nas
Partes~\ref{parte:fundamentos}--\ref{parte:mare} e o coeficiente de arrasto
linear na Parte~\ref{parte:melhorias}; $f$ é o parâmetro de Coriolis no modelo
dinâmico e o fator nodal na análise harmônica ($f_i$); $h$ é a elevação da bacia
no Capítulo~\ref{cap:oscilador}, a longitude média do Sol na
Tabela~\ref{tab:doodson} e o número de Love $h_2$ em \eqref{eq:love}. O ângulo
horário recebe símbolo próprio, $\mathcal{H}$, justamente para não colidir com a
profundidade $H$.

% ---------------------------------------------------------------------
% =====================================================================
\begin{thebibliography}{99}
\addcontentsline{toc}{chapter}{\bibname}

\bibitem{thornton2004} J.~B. Marion and S.~T. Thornton,
\emph{Classical Dynamics of Particles and Systems}, \S5.5 ``Ocean Tides'',
pp.~198--204.

\bibitem{accad1978} Y.~Accad and C.~L. Pekeris,
``Solution of the tidal equations for the M\textsubscript{2} and
S\textsubscript{2} tides in the world oceans from a knowledge of the tidal
potential alone'', \emph{Phil.\ Trans.\ R.\ Soc.\ Lond.\ A} \textbf{290}, 1978.

\bibitem{arakawa1977} A.~Arakawa and V.~R. Lamb,
``Computational design of the basic dynamical processes of the UCLA general
circulation model'', \emph{Methods in Computational Physics} \textbf{17}, 1977.

\bibitem{cartwright1971} D.~E. Cartwright and A.~C. Edden,
``Corrected tables of tidal harmonics'',
\emph{Geophys.\ J.\ R.\ Astron.\ Soc.} \textbf{33}, 1973.

\bibitem{doodson1921} A.~T. Doodson,
``The harmonic development of the tide-generating potential'',
\emph{Proc.\ R.\ Soc.\ Lond.\ A} \textbf{100}, 1921.

\bibitem{egbert2000} G.~D. Egbert and R.~D. Ray,
``Significant dissipation of tidal energy in the deep ocean inferred from
satellite altimeter data'', \emph{Nature} \textbf{405}, 775--778, 2000.

\bibitem{garrett1972} C.~Garrett,
``Tidal resonance in the Bay of Fundy and Gulf of Maine'',
\emph{Nature} \textbf{238}, 441--443, 1972.

\bibitem{gill1982} A.~E. Gill,
\emph{Atmosphere--Ocean Dynamics}, Academic Press, 1982.

\bibitem{hairer2006} E.~Hairer, C.~Lubich and G.~Wanner,
\emph{Geometric Numerical Integration}, 2nd ed., Springer, 2006.

\bibitem{jayne2001} S.~R. Jayne and L.~C. St.~Laurent,
``Parameterizing tidal dissipation over rough topography'',
\emph{Geophys.\ Res.\ Lett.} \textbf{28}, 2001.

\bibitem{lyard2021} F.~H. Lyard \emph{et al.},
``FES2014 global ocean tide atlas: design and performance'',
\emph{Ocean Science} \textbf{17}, 615--649, 2021.

\bibitem{munk1998} W.~Munk and C.~Wunsch,
``Abyssal recipes II: energetics of tidal and wind mixing'',
\emph{Deep-Sea Research I} \textbf{45}, 1977--2010, 1998.

\bibitem{pedlosky1987} J.~Pedlosky,
\emph{Geophysical Fluid Dynamics}, 2nd ed., Springer, 1987.

\bibitem{ray1998} R.~D. Ray,
``Ocean self-attraction and loading in numerical tidal models'',
\emph{Marine Geodesy} \textbf{21}, 181--192, 1998.

\bibitem{schureman1958} P.~Schureman,
\emph{Manual of Harmonic Analysis and Prediction of Tides},
U.S. Coast and Geodetic Survey, Special Publication 98, 1958.

\bibitem{stammer2014} D.~Stammer, R.~D. Ray, O.~B. Andersen, B.~K. Arbic
\emph{et al.},
``Accuracy assessment of global barotropic ocean tide models'',
\emph{Reviews of Geophysics} \textbf{52}, 243--282, 2014.

\bibitem{taylor1919} G.~I. Taylor,
``Tidal friction in the Irish Sea'',
\emph{Phil.\ Trans.\ R.\ Soc.\ Lond.\ A} \textbf{220}, 1919.

\bibitem{taylor1922} G.~I. Taylor,
``Tidal oscillations in gulfs and rectangular basins'',
\emph{Proc.\ London Math.\ Soc.} \textbf{20}, 148--181, 1922.

\bibitem{yoshida1990} H.~Yoshida,
``Construction of higher order symplectic integrators'',
\emph{Physics Letters A} \textbf{150}, 262--268, 1990.

\bibitem{etopo2022} NOAA National Centers for Environmental Information,
\emph{ETOPO 2022 15 Arc-Second Global Relief Model}, 2022.
\url{https://www.ncei.noaa.gov/products/etopo-global-relief-model}

\bibitem{noaa} NOAA CO-OPS,
\emph{Harmonic constituents and tide predictions API}.
\url{https://api.tidesandcurrents.noaa.gov/}

\end{thebibliography}

\end{document}
